\documentclass[output=paper,colorlinks,citecolor=brown]{langscibook}
\author{Michael Daniel\affiliation{Collegium de Lyon, University of Lyon} and Nina Dobrushina\affiliation{Laboratoire Dynamique du Langage (CNRS), Lyon}}
\title{Apprehensives in East Caucasian}
\abstract{In this paper, we describe the category of apprehensive in three languages of the East Caucasian (Nakh-Daghestanian) language family, Mehweb (Dargwa), Archi (Lezgic) and Akhvakh (Andic). The fact that the apprehensive, a cross-linguistically rare category, is of independent origin in these three related languages, may point to areal diffusion, and yet, the languages are not in direct contact, and the apprehensive has not yet been attested in other languages of the family. Further, in Mehweb and Archi, the two apprehensives show some important functional similarities. We attempt to reconstruct their morphosyntactic evolution, and argue that, after having followed very different paths of grammaticalization, the two apprehensives have stabilized in a similar functional domain.}



\IfFileExists{../localcommands.tex}{%hack to check whether this is being compiled as part of a collection or standalone
 \usepackage{tabularx,multicol}
\usepackage{url}
\urlstyle{same}

\usepackage{langsci-optional}
\usepackage{langsci-lgr}
\usepackage{langsci-gb4e}

\babelprovide[import]{hebrew}

\usepackage{paralist} %%may not be needed, currently used in questionnaire


%Siena added packages for introduction; not sure they are needed
\usepackage[normalem]{ulem}
\useunder{\uline}{\ul}{}

%from Treis:
\usepackage{listings}
\lstset{basicstyle=\ttfamily,tabsize=2,breaklines=true}

%%from Francez
% % % \usepackage{cjhebrew}
%\usepackage{tipa}  Also included by D&A, need to check with LSP whether to change it to language science tipa
\usepackage{leipzig}
%%from Dabkowski&AnderBois:

\usepackage[shortlabels]{enumitem} 
\usepackage{centernot}
\usepackage{arydshln}
\usepackage{newunicodechar}
\usepackage{csquotes}
%\let\sups\relax \usepackage{tipa} commented out on 7/3/25 Martina
% \usepackage[all]{nowidow}
\usepackage{listings} \lstset{basicstyle=\ttfamily}
\usepackage{multirow}
\usepackage{rotating}
\usepackage{pdflscape}
\usepackage{xspace}
%\usepackage{tabularx}
%\usepackage{multicol}
\usepackage{supertabular}
\usepackage{hhline}
\usepackage{bbding} %this package is for the checkmarks in Table 6 in Browne et al.

%\usepackage{qtree}
%\usepackage{tree-dvips}

\usepackage{array}
\usetikzlibrary{calc,tikzmark,matrix}
\usepgflibrary{shadings}
\usepackage{pifont}
%from Fedotov:

\usepackage{langsci-textipa}
%\usepackage{multirow}


%\usepackage{xeCJK}

%%%%%%%%%%%%%%%%%%%%%%%%%%%%%%%%%%%%%%%%%%%%%%%%%%%%
%%% EXAMPLES %%%%%%%%%%%%%%%%%%%%%%%%%%%%%%%%%%%%%%%
%%%%%%%%%%%%%%%%%%%%%%%%%%%%%%%%%%%%%%%%%%%%%%%%%%%% 

% if you want the source line of examples to be in
% italics, uncomment the following line
\renewcommand{\exfont}{\itshape}
% to add additional information to the right of
% examples, uncomment the following line

%\usepackage{langsci/styles/jambox}


%%% FOOTNOTE EXAMPLE NUMBERING
% \makeatletter
% \patchcmd{\@footnotetext}{\setcounter{fnx}{0}}{\renewcommand{\thexnumi}{\roman{xnumi}}}{}{}
% \apptocmd{\@footnotetext}{
%     \@noftnotetrue
%     \renewcommand{\thexnumi}{\arabic{xnumi}}%
% }{}{}
%
% \makeatother
\usepackage[linguistics,edges]{forest}
\usepackage{hyphenat}
\usepackage{langsci-branding}

  \makeatletter
\let\thetitle\@title
\let\theauthor\@author
\makeatother

\newcommand{\togglepaper}[1][0]{
   \bibliography{../localbibliography}
   \papernote{\scriptsize\normalfont
     \theauthor.
     \titleTemp.
    To appear in:
    Martina Faller,  Marine Vuillermet \&  Eva Schultze-Berndt (eds.). Apprehensional Constructions in a cross-linguistic perspective. Berlin: Language Science Press. [preliminary page numbering]
   }
   \pagenumbering{roman}
   \setcounter{chapter}{#1}
   \addtocounter{chapter}{-1}
}

\newbool{bookcompile}
\booltrue{bookcompile}
\newcommand{\bookorchapter}[2]{\ifbool{bookcompile}{#1}{#2}}


\newfontfamily\FedotovBoldEmptySetFont{LibertinusSerif-Italic.otf}[AutoFakeBold=2.5]
\newcommand{\FedotovBoldEmpty}{{\FedotovBoldEmptySetFont\bfseries ∅}}

%From Dabkowski&AnderBois


\let\sc\scshape
\let\rm\rmfamily
%\let\it\itshape
\let\tt\ttfamily
\let\nr\normalfont

\newcommand{\wsc}[1]{\textls[80]{\scshape#1}}

\let\textai\textit
% \let\textlc\spacedlowsmallcaps
% \let\textac\spacedallcaps
\let\textlc\textsc
\let\textac\textsc


%%%%%%%%%%%%%%%%%%%%%%%%%%%%%%%%%%%%%%%%%%%%%%%%%%%%
%%% FLAG %%%%%%%%%%%%%%%%%%%%%%%%%%%%%%%%%%%%%%%%%%%
%%%%%%%%%%%%%%%%%%%%%%%%%%%%%%%%%%%%%%%%%%%%%%%%%%%%

% \newcommand{\scott}  [1]{\textcolor{violet}{#1}}
% \newcommand{\maks}   [1]{\textcolor{blue}{#1}}
% % \newcommand{\data}   [1]{\textcolor{teal}{#1}}
% \newcommand{\new}   [1]{#1}
% \newcommand{\newnew}   [1]{\textcolor{teal}{#1}}
% \newcommand{\rev}   [1]{\textcolor{magenta}{#1}}



%%%%%%%%%%%%%%%%%%%%%%%%%%%%%%%%%%%%%%%%%%%%%%%%%%%%
%%% TABLES %%%%%%%%%%%%%%%%%%%%%%%%%%%%%%%%%%%%%%%%%
%%%%%%%%%%%%%%%%%%%%%%%%%%%%%%%%%%%%%%%%%%%%%%%%%%%%

%\newcolumntype{Y}{>{\arraybackslash}p{25.3543464286pt}} % length of L

% \newcolumntype{K}{>{\arraybackslash}p{24.25pt}}
% \newcolumntype{V}{>{\arraybackslash}p{(\textwidth/9-24pt)/2}}
%
\let\ch\relax \newcommand{\ch}[2]{\multicolumn{#1}{c}{\sc#2}}
\let\sx\relax \newcommand{\sx}[2]{\multicolumn{#1}{l}{\emph{-#2}}} % suffix
\let\su\relax \newcommand{\su}{\sx{1}} % suffix unicolumnal
\let\sd\relax \newcommand{\sd}{\sx{2}} % suffix dicolumnal
\let\cx\relax \newcommand{\cx}[2]{\multicolumn{#1}{l}{\emph{=#2}}} % clitic
\let\cu\relax \newcommand{\cu}{\cx{1}} % clitic unicolumnal
\let\cd\relax \newcommand{\cd}{\cx{2}} % suffix dicolumnal
\let\gx\relax \newcommand{\gx}[2]{\multicolumn{#1}{l}{\sc #2}} % gloss
\let\gd\relax \newcommand{\gd}{\gx{2}} % gloss dicolumnal
\let\gr\relax \newcommand{\gr}[1]{\multicolumn{2}{c}{\multirow{2}{*}{\textcolor{gray}{\sc#1}}}}
\let\db\relax \newcommand{\db}[1]{\multicolumn{2}{c}{\multirow{2}{*}{\sc#1}}}



%%%%%%%%%%%%%%%%%%%%%%%%%%%%%%%%%%%%%%%%%%%%%%%%%%%%
%%% EXAMPLES %%%%%%%%%%%%%%%%%%%%%%%%%%%%%%%%%%%%%%%
%%%%%%%%%%%%%%%%%%%%%%%%%%%%%%%%%%%%%%%%%%%%%%%%%%%%

%%% AFFIX BOUNDARY %%%%%%
% \DeclareRobustCommand{\longmn}{\textminus}
% \DeclareRobustCommand{\mn}{-}

%%% CLITIC BOUNDARY %%%%%%
% \DeclareRobustCommand{\longeq}{=}
% \makeatletter\DeclareRobustCommand{\eq}{%
%     \settowidth{\@tempdima}{-}% width of a hyphen
%     \resizebox{\@tempdima}{\height}{=}%
% }\makeatother

%%% FUZZY CLITIC %%%%%%
% \DeclareRobustCommand{\longfc}{%
%     $\overset{\text{\smash{}}}{\text{$\approx$}}$%
% }
% \makeatletter\DeclareRobustCommand{\fc}{%
%     \settowidth{\@tempdima}{-}% width of a hyphen
%     \resizebox{\@tempdima}{\height}{%
%         $\overset{\text{\smash{}}}{\text{$\approx$}}$%
%     }%
% }\makeatother

%%% REDUPLICATION %%%%%%%
% \DeclareRobustCommand{\longtl}{$\sim$}
% \makeatletter\DeclareRobustCommand{\tl}{%
%     \settowidth{\@tempdima}{-}% width of a hyphen
%     \resizebox{\@tempdima}{\height}{$\sim$}%
% }\makeatother

% \newcommand{\mn}          {\textminus}
\newcommand{\src}    [1]  {\jambox[0pt]{\rm(#1)}}
% \newcommand{\ai}     [2]  {\emph{#1} `\textsc{#2}#3'\xspace}
\newcommand{\ai}     [2]  {\emph{#1} `\textsc{#2}'\xspace}
\newcommand{\sane}   [1][]{\ai{=sa'ne}{appr}{#1}}
\newcommand{\srcn}   [1][]{\ai{kûintsû}{srcn}{#1}}
% \newcommand{\mbekane}[1][]{\ai{\eq mbe kañe}{neg.adv aux.inf}{#1}}

% \let\sc\relax \newcommand{\sc}{\normalfont\scshape}
% \let\rm\relax \newcommand{\rm}{\normalfont\rmfamily}
% \let\it\relax \newcommand{\it}{\normalfont\itfamily}

\newcommand{\lb}{{\rm[}}
\newcommand{\rb}{{\rm]}}

\newcommand{\bad}{\makebox[0pt][r]{*\,}}
\newcommand{\infel}{\makebox[0pt][r]{\#\,}}


% \newcommand{\lsqz}[1]{#1}
% \newcommand{\lsqz}[1]{\makebox[0pt][l]{#1}}

% \newcommand{\lsqu}[1]{\makebox[0pt][l]{#1}}
% \newcommand{\rsqu}[1]{\makebox[0pt][r]{#1}}
% \newcommand{\astr}{\makebox[0pt][r]{{\nr*}}}
% \newcommand{\hash}{\makebox[0pt][r]{{\nr\#}}}
% \newcommand{\qstn}{\makebox[0pt][r]{{\nr?}}}
% \newcommand{\bad}{\makebox[0pt][r]{*}}
% \newcommand{\unhappy}{\makebox[0pt][r]{\#}}


%%%%%%%%%%%%%%%%%%%%%%%%%%%%%%%%%%%%%%%%%%%%%%%%%%%%
%%% UNICODE CHARACTERS  %%%%%%%%%%%%%%%%%%%%%%%%%%%%
%%%%%%%%%%%%%%%%%%%%%%%%%%%%%%%%%%%%%%%%%%%%%%%%%%%%

\newunicodechar{⟨}{⟨}
\newunicodechar{⟩}{⟩}

\newcommand{\infix}{⟨}
\newcommand{\outfix}{⟩}



%%%%%%%%%%%%%%%%%%%%%%%%%%%%%%%%%%%%%%%%%%%%%%%%%%%%
%%% IPA %%%%%%%%%%%%%%%%%%%%%%%%%%%%%%%%%%%%%%%%%%%%
%%%%%%%%%%%%%%%%%%%%%%%%%%%%%%%%%%%%%%%%%%%%%%%%%%%%

% \newcommand{\ipa}{\textipa}

\newcommand{\sub}{\textsubscript}
% \let\super\relax \newcommand{\super}{\textsuperscript}

\newcommand{\ts}{\texttslig}
% \let\dz\relax \newcommand{\dz}{\textdzlig}
\newcommand{\tS}{\textteshlig}
\newcommand{\dZ}{\textdyoghlig}
\newcommand{\cj}{\textctc}
\newcommand{\sj}{\textctc}
\newcommand{\zj}{\textctz}
\newcommand{\tcj}{\texttctclig}
\newcommand{\tsj}{\texttctclig}
\newcommand{\dzj}{\textdctzlig}
\newcommand{\gj}{\textbardotlessj}
% \let\nj\relax \newcommand{\nj}{\textltailn}
% \newcommand{\W}{\textturnmrleg}
\newcommand{\wl}{\textturnmrleg}
% \newcommand{\1}{\ipabar{\~\i}{.5ex}{1.1}{}{}}
\newcommand{\dotlessibar}{\ipabar{\i}{.5ex}{1.1}{}{}}
\newcommand{\darkl}{\textltilde}
% \newcommand{\n}{\textsubarch}
% \newcommand{\x}{\textovercross}
\newcommand{\lowered}{\textlowering}
\newcommand{\raised}{\textraising}
\newcommand{\retracted}{\textsubbar}
\newcommand{\unreleased}{\textcorner}
% \newcommand{\syllabic}{\s}
\newcommand{\implosivet}{\texthtt}
\newcommand{\implosived}{\texthtd}
\newcommand{\implosivep}{\texthtp}
\newcommand{\implosiveb}{\texthtb}
\newcommand{\implosivek}{\texthtk}
\newcommand{\implosiveg}{\texthtg}



%%%%%%%%%%%%%%%%%%%%%%%%%%%%%%%%%%%%%%%%%%%%%%%%%%%%
%%% OTHER %%%%%%%%%%%%%%%%%%%%%%%%%%%%%%
%%%%%%%%%%%%%%%%%%%%%%%%%%%%%%%%%%%%%%%%%%%%%%%%%%%%

\newcommand{\ie}{i.\,e.\xspace}
\newcommand{\Ie}{I.\,e.\xspace}
\newcommand{\eg}{e.\,g.\xspace}
\newcommand{\Eg}{E.\,g.\xspace} 




\newcommand{\francoissource}[1]{\hfill\textnormal{{\footnotesize #1} }}
\newcommand{\E}{Ɛ\kern-1.5pt\raisebox{1pt}{̏}\kern1.5pt}
\newcommand{\ù}{{ṵ}\kern-2pt{̀}\kern2pt}
\newcommand{\ȕ}{{ṵ}\kern-2pt{̏}\kern2pt}

\newcommand{\OO}{ɔ\kern-1pt\raisebox{-1pt}{̰}\kern1pt}
\newcommand{\Ǒ}{{{ɔ̌\kern-1pt\raisebox{-.5pt}{̰}\kern1pt}}}


\newcommand{\à}{{a̰}\kern-2pt{̀}\kern2pt}
\newcommand{\ilit}[1]{#1\il{#1}}
  \togglepaper[23]
 }{}

\begin{document}
\maketitle
%\shorttitlerunninghead{}%%use this for an abridged title in the page headers


\section{Introduction}

 
Apprehensives are usually described as a morphological expression of fear (\citealt{lichtenberk1995,Dobrushina2006cyrillic,vuillermet2018a,Smith-Dennis2021}). With an apprehensive, the speaker reports that there is a probability of an undesirable event and that he or she fears it. Apprehensives may simultaneously convey the speaker's emotion (they are apprehensive of a potentially negative event) and a command (they induce the addressee to take measures to avoid or prevent the undesirable situation). Due to the directive component (expression of a command), apprehensives may be part of main clause phenomena. Cross-linguistic variation in their status as forms used exclusively in the main clause or also in subordination is particularly relevant for their typology. Dedicated morphological apprehensives are cross-linguistically infrequent. 

This paper discusses dedicated apprehensives in three languages of the East Caucasian language family: Archi (Lezgic branch of the family), Mehweb (Dargwa branch), and North Akhvakh (Avar-Andic branch). Although the three languages belong to the same family and roughly the same geographic region (central to northern highlands of Daghestan, Eastern Caucasus), they are only distantly related -- belonging to different branches of this relatively old family -- and are not spoken in geographically adjacent areas. All three of them are minority languages. Archi and Mehweb are one-village languages, and Northern Akhvakh is spoken in only four villages (but see below). According to the 1926 Census, which reflects traditional population sizes better than later statistics, Mehweb was spoken by 710, Archi by 856, and Northern Akhvakh by 2216 people (\citealt{materials1927}, \citealt{Dobrushinaetal2017}). Today, there are many more villagers living in the lowlands, but the degree of language vitality among resettlers is unclear. So far, these languages are the only languages of the family where inflectional apprehensives have been attested, but this might be due simply to the fact that apprehensives are often overlooked by descriptive grammars.

In Archi, Mehweb and Akhvakh, the apprehensives are expressed morphologically by verbal suffixes \nobreakdash-\textit{lkːut}, \nobreakdash-\textit{ala} and \nobreakdash-\textit{alago}(\textit{le}), respectively. Based on our field data as well as, to a certain extent, on the existing descriptive sources, we will describe the uses typical of these forms, review their availability with subjects of different persons, consider whether they are compatible with negation, discuss their status as finite or non-finite forms, and put forward hypotheses about the origins of the markers. In spite of the apparent phonological similarity of the three affixes, there is no evidence that they are cognate. We will argue that they result from three different paths of grammaticalization from different -- even if not always known -- sources. In Archi, the independent apprehensive probably arose through insubordination (expansion from syntactically dependent to main clause uses -- \citealt{evans2007insubordination}). In Mehweb, on the other hand, subordinate uses of the apprehensive seem to evolve through an extension of the reported speech construction. In Akhvakh, its probable origin is univerbation of a verb with an interjection expressing warning. At least in Archi and Mehweb, these forms, although coming from very different sources, ultimately stabilize in very similar functional domains typical of the apprehensive cross-linguistically. Presently, it remains unclear how much of areal impact this process may have had.

For Mehweb, corpus data are limited due to the corpus size (own data, \citealt{Magometov1982}), so we mostly illustrate the uses of the apprehensive with elicited data. For Northern Akhvakh, no corpus was available, but elicited data are supplemented by several examples from a dictionary \citep{MagomedovaAbdulaeva2007}. For Archi, we have the by far richest attestation of apprehensives, so that, where available, corpus examples are provided together with elicited data, but even here naturalistic examples are limited to 26 contexts. Archi corpus data integrate various sources, including texts published in \citet{Dirr1908}, \citet{Mikailov1967}, \citet{Kibrik1977}, and \citet{SamedovMagdilova2017a, SamedovMagdilova2017b} as well as our own unpublished data; all these texts will be collectively referred to as \textit{corpus}, and the examples coming from them are marked as \textit{naturalistic} as opposed to \textit{elicited}.

\sectref{sec:Archi},  \sectref{sec:Mehweb}, and \sectref{sec:NAkhvakh} describe Archi, Mehweb and Northern Akhvakh data, respectively. Note that the data were collected at very different times, and thus not all types of relevant information available from one language is also available from another. \sectref{daniel:summary} is a comparative discussion of the data of the three languages, and \sectref{daniel:conclusion} is a brief conclusion.




\section{Archi: From `fear' clauses to apprehensive via insubordination}
\label{sec:Archi}
 
The Archi form in \textit{-lkːut} was described by \citet[292]{Kibrik1977} as a converb used in negative purpose constructions. While Kibrik only quotes subordinate uses (see \sectref{subsec:Asubordinate} below), \citet{Dobrushina2006cyrillic} shows that, when used independently, this form is interpreted as apprehensive:\footnote{Here and below, for Archi, numbers in the glosses refer to genders. Typically, Gender 1 is for male referents, and Gender 2 is for female referents, and the rest of nouns are distributed between Gender 3 and 4. Mehweb and Northern Akhvakh only have three genders, designated as M, F and N. In all three languages, HPL stands for human plural agreement which conflates male and female referents, and NPL stands for non-human plural agreement.}


\ea \label{ex:DD1} \textnormal{elicited} \\
\gll jami č'abu \textbf{esa-lkːut}! \\
wolf.\textsc{erg} sheep.\textsc{pl} \textsc{npl}.slaughter-\textsc{appr} \\
\glt `The wolf might kill the sheep!' \\
\z

The origin of the form in \textit{-lkːut} is unclear. The suffix is added to what Kibrik sees as the infinitive stem of the verb (most other forms are based on the perfective or the imperfective stem). It is very tempting to suggest univerbation of the lexical verb with the verb `see' (perfective stem \textsc{CL}-\textit{ak:u}) with loss of the class prefix and initial vowel. First of all, there is a cross-linguistically robust connection between `see' and apprehensives \citep{ lichtenberk1995,Dobrushina2006cyrillic}. There is another univerbation path that `see' follows in Archi, which leads to a series of \textit{verificative} forms meaning `find out whether  \ldots' (see \citealt{Maisak2016} on the category of verificative in Aghul and several other East Caucasian languages, also \citealt{Kibrik1977}). This may be taken as a token of the propensity of the verb `see' to grammaticalize. This stipulation however remains problematic because of the status of the additional material (initial \textit{l}- and the final \nobreakdash-\textit{t}) which remains unclear. In fact, in view of our tentative conclusions on the evolution of this form below in \sectref{subsec:Aapprehensive}, there might be no need to look for the source of the apprehensive suffix among the known lexical sources of its grammaticalization in other languages, because its development towards apprehensive uses might have occurred later, through an evolution of a syntactic construction.   

In Archi, the apprehensive is used both in main clauses and in some subordinate clauses, including negative purpose clauses, `in-case' clauses (as termed in \citealt{lichtenberk1995}) and the complement clauses of verbs of `fear'. Except `in-case' clauses, all types are attested in the corpus; but, as we will see, it is not always easy to distinguish between these types, and all subordinate uses may be argued to fulfill the same general function.



\subsection{Semantics and morphology}
\label{subsec:Asemmorph}

In independent uses, the Archi form in \textit{-lkːut} expresses the fact that the speaker is anticipating a negative situation. Out of some 20 occurrences of the apprehensive in the corpus, only one is arguably independent, while all other examples come from subordinate clauses. We do not believe this reflects anything except the nature of our corpus, which is almost exclusively narrative, and not interactional. The naturalistic example of independent use comes embedded in a speech reporting construction. 


\ea \label{ex:DD2} \textnormal{naturalistic} \\
\gll χir q'ˁʷebos priq'at e{\textlangle}r{\textrangle}tːi jasːa, χir wa buwa, un=ur=u χːal-ši \textbf{e{\textlangle}r{\textrangle}χːu-lkːut}, t'o=ra=r, un=ur=u χːal-ši \textbf{e{\textlangle}r{\textrangle}χːu-lkːut}, t'o=ra=r  \\
then next calm {\textlangle}2{\textrangle}become:\textsc{pfv} now then oh Mom you.sg\textsc{=quot=add} be.bad-\textsc{cvb} \textbf{{\textlangle}2{\textrangle}remain\textsc{-appr}} no=\textsc{q=quot} you.sg\textsc{=quot=add} be.bad-\textsc{cvb} \textbf{{\textlangle}2{\textrangle}remain\textsc{-appr}} no=\textsc{q=quot} \\
\glt `Then she calmed down, then Oh Dear! (lit.\ `Oh mother!'), she said, may you not feel bad, right (lit.\ `is it not')? may you not feel bad, right?'
\z

 Example \REF{ex:DD1} shows the use of apprehensive in the third person, and \REF{ex:DD2} in the second person. In elicitation, the apprehensive easily combines with subjects of all persons:
 
\ea \label{ex:DD3} \textnormal{1st person, elicited} \\
\gll bara zon \textbf{d-ukːu-lkːut}! \\
beware I:\textsc{nom} \textbf{2-burn-\textsc{appr}} \\
\glt `Careful, I might get burnt!' (when another person carries boiling water)
\z

\ea \label{ex:DD4} \textnormal{2nd person, elicited} \\
\gll un oˁrču-tːu-t ɬːan c'abu, bara un bec'ːo-ši \textbf{kʷe-lkːut} \\
you.sg\textsc{(erg)} 4.freeze:\textsc{pfv-attr-4} water drink:\textsc{pfv} beware you.sg\textsc{.nom} get.sick-\textsc{cvb} \textbf{1.become-\textsc{appr}} \\
\glt `You drank cold water, you might get ill.'
\z

\ea \label{ex:DD5} \textnormal{3rd person, elicited}\\
\gll jaˁt'i-li \textbf{w-eq'ːi-lkːut}! \\
snake-\textsc{obl.erg} \textbf{1-bite-\textsc{appr}} \\
\glt `The snake might bite (me)!'
\z


 With the 2nd person, the apprehensive \textit{iwkulkːut} `Watch out not to fall!' in \REF{ex:DD6} is close to the prohibitive \textit{werkurdigi} `Do not fall!' in that the apprehensive may be taken as a directive to avoid a dangerous outcome. For a comparison of the apprehensive and the prohibitive, it is not easy to obtain speakers' explicit intuitions in a comparable form. We asked one of our consultants which of the forms implied a ``higher danger''. For her, it was the apprehensive, while the prohibitive, as she had put it, suggests that even if the addressee falls, she would simply get up again, end of story (Russian \textit{упадет и встанет}). Another consultant commented that, when using the prohibitive, the speaker is ``sure that the addressee will fall down unless a precaution is taken'', while in the case of the apprehensive, she ``merely suggests it might happen''. Finally, yet another consultant, when asked to compare the two forms, simply suggested the examples in \REF{ex:DD6} as natural contexts:



\ea \label{ex:DD6} \textnormal{examples suggested by the consultant} \\
\ea \label{ex:DD6a}
\gll \textbf{i{\textlangle}w{\textrangle}ku-lkːut} qʷat'i-li-tːi-k w-e{\textlangle}r{\textrangle}χi-r-digi \\
\textbf{{\textlangle}1{\textrangle}fall-\textsc{appr}} tree-\textsc{obl-sup-lat} 1-{\textlangle}\textsc{ipfv}{\textrangle}climb-\textsc{ipfv-proh} \\
\glt `Don't climb up the tree, lest you fall down.' \\
\ex \label{ex:DD6b}
\gll doːˁzu-t minnat i, \textbf{ w-e{\textlangle}r{\textrangle}ku-r-digi}! \\
big-4 request 4.be \textbf{ 1-{\textlangle}\textsc{ipfv}{\textrangle}fall-\textsc{ipfv-proh}} \\
\glt `Let me ask you one thing -- do not fall down!'
\z
\z


 These intuitions and suggestions may be linked to the epistemic and illocutionary components present in apprehensives vs.\ prohibitives. Apprehensives inform of a probability of a negative event in the future, while prohibitives, like imperatives, are more immediately connected to the moment of speech and, in these cases, seem to be used to prevent a situation. It thus seems that the apprehensive form conveys the feeling of danger in a more direct way (more of a warning), while the prohibitive is an instruction to avoid the situation (more of a command), but we admit that this difference can at least partly be in the eye of the beholder. 

One of the reviewers remarks that a natural expectation of the difference between the apprehensive and the prohibitive is that of exerting control over the situation, because, as any directive, the prohibitive implies control. Indeed, even though the verb in \REF{ex:DD6a} and \REF{ex:DD6b} is the same, one may suspect that \REF{ex:DD6b} represents the situation of falling or not falling as more controllable, and the effect is probably obtained exactly by forming a prohibitive of a verb that normally describes a situation beyond control (a similar effect would probably be obtained in some languages by forming an imperative `Fall!' as a warning of an explosion).  Presence or absence of the subject's control may thus indeed be involved; however note that, in Archi, the prohibitive is compatible with uncontrollable states, as in \REF{ex:DD7}. To our eyes, speculating that this context also implies control (e.g.\ \textit{have fun! do something so that you are not bored!}) leaves room to extend the concept of control to any situation at all.


\ea \label{ex:DD7} \textnormal{naturalistic}\\
\gll os t'i-tːu-t sːaʕal-li-t \textbf{bazar} \textbf{kʷe-r-digi} zon os pulan-nu-t biq'ʷ-ma-k uqˁa-na, u{\textlangle}w{\textrangle}\textsc{l}i-lkan\\
one small-\textsc{attr}-4 time-\textsc{obl-sup} \textbf{bored} \textbf{1.become-\textsc{ipfv-proh}} I:\textsc{nom} one definite-\textsc{attr}-4 place-\textsc{loc-lat} 1.go:\textsc{pfv-cvb} {\textlangle}1{\textrangle}come-\textsc{temp}\\
\glt `Take care of yourself (lit.\ `do not become bored'), while I go somewhere else and come back.'
\z

 In Archi, apprehensives can refer to situations that cannot be controlled, such as meteorological events in third person constructions:

  \begin{exe}\ex \label{ex:DD8} \textnormal{elicited}\\
{\gll mač'a \textbf{ke-lkːut,} nen-t'u no\textsc{l}'-i-ši χatːi-tːa\\
be.dark 4.become-\textsc{appr} we-\textsc{ptcl.4.incl} house-\textsc{in-all} \textsc{locpl}.go.\textsc{fut-temp}\\\glt `It might become dark when we go home.'}
\end{exe}


 Most verbal forms have an inflectional negative polarity counterpart, but there is no negative apprehensive. To warn about a risk that something desirable might not happen, Archi uses a periphrastic construction of a negative converb of the main verb with an apprehensive form of the verb `remain'.\footnote{One of the reviewers asks whether the construction is monoclausal periphrastic or biclausal. Though we have no relevant syntactic evidence, we observe that in the few naturalistic examples we have, `remain' strictly follows the negative converb of what we assume to be the lexical verb in the periphrastic construction; and that the meaning of the construction (`remains not having V-ed', e.g.\  lit.\ `lest she remains not having returned' in \REF{ex:DD10}) is not, or is not fully, compositional. Before more syntactic data are obtained, this may be taken as evidence in favor of a monoclausal interpretation.} Cf.\ \REF{ex:DD9} with a regular apprehensive and \REF{ex:DD10} with a palliative `negative apprehensive', and also \REF{ex:DD11}:

\ea \label{ex:DD9} \textnormal{elicited} \\
\gll t'eˁ-tːu caqˁ-a, χːˁel \textbf{eχmu-lkːut} \\
flower-\textsc{pl} \textsc{npl}.cover-\textsc{imp} rain \textbf{4.rain-\textsc{appr}} \\
\glt `Cover the flowers, it may rain.'
\z


\ea \label{ex:DD10} \textnormal{elicited} \\
\gll t'eˁ-tː-e-s ɬːan ec-a, χːˁel \textit{eχ-di-t'uw} \textbf{eχːə-lkːut} \\
flower-\textsc{pl-obl.pl-dat} water 4.pour-\textsc{imp} rain \textbf{4.rain-\textsc{pfv-cvb.neg}} \textbf{{\textlangle}2{\textrangle}remain-\textsc{appr}} \\
\glt `Water the flowers, what if it does not rain today.' (lit.\ `lest it remains without the rain having rained.')
\z
\ea \label{ex:DD11} \textnormal{elicited} \\
\gll t'eˁ-tː-e-s ɬːan eca, to-r \textbf{boq'ˁo-t'aw} \textbf{e{\textlangle}r{\textrangle}χːə-lkːut}\\
flower-\textsc{pl-obl.pl-dat} water 4.pour-\textsc{imp} this-2 \textbf{come.back:\textsc{pfv-cvb.neg}} \textbf{4.remain-\textsc{appr}}\\
\glt `Water the flowers, what if she does not come back.'
\z 



In Archi, most forms usually considered finite may take the negative suffix \nobreakdash-\textit{t'u.} The fact that, in negative contexts, the apprehensive takes a periphrastic form with a special negative converb may be indicative of the non-finite status of the apprehensive or the form it originated from. Also note that, in this construction, the apprehensive form of the main verb seems to be synonymous to the infinitive of the same periphrastic construction (see \REF{ex:DD16} in the context of negative purpose).

As these examples show, the apprehensive is often used in combinations with clauses that motivate the usage of the apprehensive by specifying the exact situation that can lead to an undesirable outcome; see also \REF{ex:DD12}:
 
    \ea \label{ex:DD12} \textnormal{elicited} \\
\gll k'umk'um \textsc{l}'eˁr-ši, \textbf{d-ok:i-lkːut} \\
pan be.hot-\textsc{ccvb.4.aux} \textbf{2-burn-\textsc{appr}} \\
\glt `The pan is hot, you might burn yourself.'
\z


 Alternatively, the non-apprehensive part can describe a way to avoid the undesirable outcome. The constructions in \REF{ex:DD9} to \REF{ex:DD12} are close to various types discussed in the next section, and at times, their status as paratactic juxtaposition (as described in this section) or subordination (as described in the following section) is unclear.
 

\subsection{Subordinate uses}
\label{subsec:Asubordinate}
 
Crosslinguistically, apprehensives are commonly used in subordination. In such contexts, the apprehensive is essentially a way to express the reason of the situation introduced by the main clause, which may describe an action undertaken to avoid the negative situation (negative purpose clauses) or, if it cannot be avoided, to reduce its negative impact (\citegen{lichtenberk1995} `in-case' clauses), or to simply describe the emotional state (some but not all `fear' clauses). In Archi, all such uses of the apprehensive are attested. In these contexts, the apprehensive may refer to the past, so that their interpretation as directives is clearly excluded. Cf.\ the naturalistic example in \REF{ex:DD13}:

\ea \label{ex:DD13} \textnormal{naturalistic} \\
\gll χːams že-n-t'u lo=wu uɬdu=wu barha-s šːʷi-šːʷi boχːˤo b-o{\textlangle}r{\textrangle}qˤi-r-ši e-b-di-li qi-qi jamu \textbf{ɬːummu-lkːut} sːak'u q'a{\textlangle}b{\textrangle}q'i-r-ši e-b-di-li\\
bear self.\textsc{obl-gen-ptcl4} child=\textsc{add} shepherd=\textsc{add} care-\textsc{inf} by.night-by.night hunt 3-{\textlangle}\textsc{ipfv}{\textrangle}go-\textsc{ipfv-cvb} be-3-\textsc{pst-evid} at.day-at.day that.1 \textbf{1.flee-\textsc{appr}} on.guard {\textlangle}3{\textrangle}sit-\textsc{ipfv-cvb} be-3-\textsc{pst-evid}\\
\glt `The she-bear, to sustain the shepherd (her husband) and child, used to go out hunting at night, and during the day she was looking after her husband \textbf{so that he might not escape}.'
\z

\noindent Example \REF{ex:DD13} illustrates the claim in \citet{vuillermet2018a}: unlike in its independent uses, in subordinate uses of the apprehensive the judge of negative evaluation may be the subject of the main clause and not the speaker herself. Indeed, the she-bear is a negative character who took the protagonist prisoner and whose evaluations of what is negative and what is positive the narrator presumably does not share. In Archi, the evaluation probably has to be done from the viewpoint of who is the referent of the subject of the main clause (independently of whether the narrator shares it or not); at least this seems to be the case in all our examples. 

In \REF{ex:DD13} above, two main clauses introduce purpose clauses, one of which is positive (`went out hunting in order to sustain') and the other is negative (`stay on guard in order that he does not escape'). Because of its built-in component of negative evaluation, the apprehensive may only be used to express negative purpose. In examples \REF{ex:DD14} and \REF{ex:DD15}, the apprehensive used as negative purpose (\REF{ex:DD14a}, \REF{ex:DD15a}) is contrasted with the use of an infinitive as its positive purpose counterpart (\REF{ex:DD14b}, \REF{ex:DD15b}):
 

\ea \label{ex:DD14} \textnormal{elicited} \\
\ea \label{ex:DD14a}
\gll buwa lo \textbf{čučəbo-lkːut} o{\textlangle}r{\textrangle}qˤa \\
mother child \textbf{wash-\textsc{appr}} {\textlangle}2{\textrangle}go:\textsc{pfv} \\
\glt `The mother went away in order not to wash the child' (lit.\ `lest she would have to wash the child'). \\
\ex \label{ex:DD14b}
\gll buwa lo \textbf{čučəbo-s} o{\textlangle}r{\textrangle}qˤa \\
mother child \textbf{wash-\textsc{inf}} {\textlangle}2{\textrangle}go:\textsc{pfv} \\
\glt `The mother went away to wash the child.'
\z
\z


\ea \label{ex:DD15} \textnormal{elicited} \\
\ea \label{ex:DD15a}
\gll buwa lo \textbf{čučəbo-lkːut} o{\textlangle}r{\textrangle}qˤa-t'u \\
mother child \textbf{wash-\textsc{appr}} {\textlangle}2{\textrangle}go:\textsc{pfv-neg}\\
\glt `The mother didn't go in order not to wash the child' (lit.\ `lest she would have to wash the child'). (i.e.\   She stayed so as not to have to wash the child.)\\
\ex \label{ex:DD15b}
\gll buwa lo \textbf{čučəbo-s} o{\textlangle}r{\textrangle}qˤa-t'u \\
mother child \textbf{wash-\textsc{inf}} {\textlangle}2{\textrangle}go:\textsc{pfv-neg}\\
\glt `The mother didn't go to wash the child.'
\z
\z

\noindent To explain the difference between the two examples in \REF{ex:DD15}, consultants indicate that, in \REF{ex:DD15a} which uses the apprehensive, the reason for the fact that the mother did not leave is clear: she did not want to deal with washing the child. In \REF{ex:DD15b}, the reason is not specified; it can be her choice not to go, but it can also be that she wanted to go and wash the child but for some reason could not make it there. Examples \REF{ex:DD13} to \REF{ex:DD15} show that the apprehensive is used in the clauses with a special type of purposive semantics, those that express situations that the subject of the main clause aims at avoiding, and the main clause describes an action undertaken to avoid it.\footnote{This clause type is probably better termed \textit{anti-purpose} rather than simply negative purpose clause, because, were this combination of propositional negation and purposive clause  compositional, it would yield something like `P, and it was not in order to Q' rather than what is usually called negative purpose, `P in order not to Q'. Similarly, prohibitives, often also called \textit{negative imperatives}, are not merely negative counterparts of the imperatives and are often referred to by a dedicated term, i.e.\   \textit{prohibitive}. In the case of the prohibitive, propositional negation simply cannot compositionally combine with imperative illocution. In fact, Kibrik considered the Archi apprehensive to be a \textit{preventive-purpose} converb.}

Another, apparently less frequent means to express the apprehensive meaning is a combination of the negative converb with the infinitive of \textit{e}${\prec}$\textsc{cl}${\succ}$\textit{χːa-} `remain', the periphrastic construction used in \REF{ex:DD10} and \REF{ex:DD11} above. Interestingly, \REF{ex:DD16} comes from the same text and describes the same situation as \REF{ex:DD26} below, where the apprehensive `fear' clause is used. Example \REF{ex:DD16} comes several clauses after the use of \REF{ex:DD26}, as a kind of an afterthought:

\ea \label{ex:DD16} \textnormal{naturalistic} \\
\gll … adam-li-s \textbf{akːu-t'aw} \textbf{eχːa-s} \\
… people-\textsc{obl-dat} \textbf{4.see:\textsc{pfv-cvb.neg}} \textbf{4.remain-\textsc{inf}} \\
\glt `… so that no one sees (the cattle).'
\z


 The situation encoded by means of the apprehensive marker is not necessarily one that instills fear in the speaker. In the following examples, \nobreakdash-\textit{lkːut} refers to a situation which is merely undesirable, not frightening:
 

\ea \label{ex:DD17} \textnormal{elicited} \\
\gll durba χːʷalli b-aχːa-qi, nen-ə{\textlangle}b{\textrangle}ij=bu. a{\textlangle}b{\textrangle}a-s \textbf{kʷa{\textlangle}b{\textrangle}šu-lkːut} \\
\textsc{hort} bread 3-take:\textsc{pfv-fut} we-{\textlangle}3{\textrangle}\textsc{ptcl.incl}={\textlangle}3{\textrangle}\textsc{ptcl} {\textlangle}3{\textrangle}do-\textsc{inf} \textbf{{\textlangle}3{\textrangle}must-\textsc{appr}} \\
\glt `Let's buy bread, so we do not have to bake it ourselves.'
\z


\ea \label{ex:DD18} \textnormal{elicited} \\
\gll durba o:k'ur herqˁu-qi, q'ˁas \textbf{ke-lkːut} nen-t'u \\
\textsc{hort} slowly \textsc{locpl}.walk:\textsc{pfv-fut} be.tired \textbf{\textsc{locpl}.become-\textsc{appr}} we-\textsc{ptcl.locpl.incl} \\
\glt `Let's go slowly, in order not to get tired.'
\z

    \ea \label{ex:DD19} \textnormal{naturalistic}\\
\gll akːon-ni-ɬːu, barq e-b-di b-akːu-t'aw, alšun-n-in hurulʕen-til b-a{\textlangle}r{\textrangle}\textsc{l}i-r-ši e-b-di-li ja-r-mi-ra-k, ja-r \textbf{bazar} \textbf{de-ke-lkːut} \\
light-\textsc{obl-comit} sun be-3-\textsc{pst} 3-see:\textsc{pfv-cvb.neg} paradise-\textsc{obl-gen} angel-\textsc{pl} \textsc{hpl}-{\textlangle}\textsc{ipfv}{\textrangle}come-\textsc{ipfv-cvb} be-\textsc{hpl-pst-evid} this-2-\textsc{obl-cont-lat} this-2 \textbf{bored} \textbf{2-become-\textsc{appr}} \\
\glt `In the morning, immediately after the sunrise, the angels of the paradise would come to her lest she felt bored.'
\z


 As a comparison of examples \REF{ex:DD14}, \REF{ex:DD15}, \REF{ex:DD17} and \REF{ex:DD18}, on the one hand, and \REF{ex:DD13} and \REF{ex:DD19}, on the other, shows, negative purpose clauses with the apprehensive may have either the same or a different subject as the main clause. More such examples (elicited data) are available in \citet[292]{Kibrik1977}.

The apprehensive may also be used in Lichtenberk's `in-case' clauses. `In-case' clauses refer to potential negative circumstances that cannot themselves be avoided, such as natural phenomena in \REF{ex:DD20}; what can be avoided or at least diminished is their negative impact. There are no clear instances of `in-case' clauses in the corpus, but they are easily obtained by elicitation, as in \REF{ex:DD20} (cf.\ also \REF{ex:DD9}--\REF{ex:DD11} above).

\ea \label{ex:DD20} \textnormal{elicited} \\
\gll č'il gʷaq'ː-a,  χːˁel \textbf{exmu-lkːut} \\
hay harvest-\textsc{imp} rain \textbf{4.rain-\textsc{appr}} \\
\glt `Collect the hay, in case it rains. (So that the hay does not get wet.)'
\z

 The apprehensive is used in dependent clauses subordinate to the verb `fear' (as in \REF{ex:DD21} and \REF{ex:DD22}) or similar verbs (as in \REF{ex:DD23}) as a matrix verb, and, judging from the meager corpus data we have, these uses are typical. In such cases, the apprehensive may provide the reason for the state of fear:

    \ea \label{ex:DD21} \textnormal{naturalistic} \\
\gll jemme-tːu-b zulmu \textbf{a{\textlangle}b{\textrangle}\textsc{l}e-lkːut} \textsc{l}'inč'a-r-š-e-b-di χalq' \\
thus-\textsc{attr}-3 violence {\textlangle}3{\textrangle}come-\textsc{appr} be.afraid-\textsc{ipfv-cvb}-be-\textsc{hpl-pst} people \\
\glt `People were afraid that such violence may happen.'
\z
\newpage
\ea \label{ex:DD22} \textnormal{elicited} \\
\gll zon \textbf{\textsc{l}'inč'a-r,} un haˁtər-čaj \textbf{də-χə-lkːut} \\
I \textbf{be.afraid}-\textsc{ipfv} you.sg\textsc{.nom} river-\textsc{obl.erg} \textbf{2-carry.away-\textsc{appr}} \\
\glt `I'm afraid that a river will carry you away.'
\z

\ea \label{ex:DD23} \textnormal{elicited} \\
\gll buwa-n \textbf{ik'ʷ} \textbf{arha-r-ši} χ:ʷak:-a lo \textbf{q'ˤʷek'e-lkːut} \\
mother-\textsc{gen} \textbf{heart} \textbf{4.think-\textsc{ipfv-cvb}.4.be} forest-\textsc{in(ess)} child \textbf{1.be.lost-\textsc{appr}} \\
\glt `The mother worried that the child might get lost in the wood.' 
\z


\noindent As in negative purpose constructions, with verbs of `fear', the form can be used in both different \REF{ex:DD21}--\REF{ex:DD23} and same \REF{ex:DD24} subject constructions:


\ea \label{ex:DD24} \textnormal{elicited} \\
\gll peč-li-ra-k hara:ši de-ke-s \textbf{\textsc{l}'inč'a-r} zon \textbf{du-k:u-lkːut} \\
stove-\textsc{obl-cont-lat} in.front 2-become-\textsc{inf} \textbf{be.afraid-\textsc{ipfv}} I(\textsc{nom}) \textbf{2-burn-\textsc{appr}} \\
\glt `I am afraid to come near the stove (because) I might get burnt.'
\z


 It is difficult to decide whether all naturalistic contexts with the verb \textit{\textsc{l}'inčar} `fear' imply a true emotion, and to what extent. Example \REF{ex:DD25} seems to unequivocally suggest a strong emotion (`I am frightened because P is likely to happen'), while example \REF{ex:DD26} may be interpreted as merely providing a reason `they did this out of fear of P'). Yet, even in \REF{ex:DD25}, one cannot be sure that the emotion is a semantic component of the apprehensive rather than simply contributed by lexical means (`trembling to death').


\ea \label{ex:DD25} \textnormal{naturalistic} \\
\gll hu, d-asːə-r-š=er inž d-ak'a-qi-ši d-i=r, χir e-r-di-na \textbf{daːɬə-de-ke-lkːut} \textbf{\textsc{l}'inč'a-r-ši} \\
yes 2-tremble-\textsc{ipfv-cvb=quot} self 2-die:\textsc{pfv-fut-cvb} 2-be=\textsc{quot} after be-2-\textsc{pst-cvb} \textbf{beat-2-\textsc{lv-appr}} \textbf{be.afraid-\textsc{ipfv-cvb}} \\
\glt `Yes, I am trembling to death, she says, afraid (i.e.\   experiencing fear?) of her getting to me and giving me a beating.'
\z

\newpage
\ea \label{ex:DD26} \textnormal{naturalistic} \\
\gll məcːor-um-čij=u ža-n-t'u buc'ːi č'abu a{\textlangle}r{\textrangle}k'u-r-tːu-t adam-li-s \textbf{akːu-lkːut} \textbf{\textsc{l}'inč'a-r-ši} \\
pasture-\textsc{pl-obl.in(ess)=add} self\textsc{.pl-gen-4.ptcl} cow.\textsc{pl} sheep.\textsc{pl} {\textlangle}\textsc{ipfv}{\textrangle}\textsc{npl}.drive-\textsc{ipfv-attr}-4 person-\textsc{obl-dat} \textbf{\textsc{npl}.see-\textsc{appr}} \textbf{be.afraid-\textsc{ipfv-cvb}}  \\
\glt `And (they) were afraid (i.e.\   thought it was possible?) that people who drove cattle in the field would see (them).'
\z


\noindent Weakening of the emotional component probably provides a bridge from the apprehensive in `fear' clauses both to the independent apprehensive and to the uses of the apprehensive in negative purpose clauses which do not involve the verb \textit{\textsc{l}'inč'ar} `fear' (such as \REF{ex:DD13} or \REF{ex:DD19}, to take only naturalistic examples). We suggest that the semantic impact of the verb of `fear' gradually weakens across constructions and uses, so that in the end the use of the verb becomes optional and the formerly subordinated apprehensive becomes insubordinated.

Indeed, in many naturalistic contexts, the verb `fear' is used non-finitely, in subordination to the main verb expressing the action taken to avoid an undesirable situation, but also subordinating an apprehensive clause that describes this situation. Syntactically, these constructions are different from those in examples \REF{ex:DD13} to \REF{ex:DD19} only in that the relation between the main and the apprehensive clause is mediated by the `fear' clause. Cf.\ the following example:

\ea \label{ex:DD27} \textnormal{naturalistic} \\
\gll  \textsc{l}'inč'a-r-ši tu-w qačaʁ=u w-i-muχur
ta-w bišin neq’ʷi-tː=u b-i-muχur, os oˁhli be-ke-\textbf{lkːut}
ža-s \textbf{he\textsc{l}'əna-lkːut} \textbf{\textsc{l}'inč'a-r-ši} 
ža-b ɬʷːak e{\textlangle}b{\textrangle}tːi-t'o=r \\
 be.afraid-\textsc{ipfv-cvb} this-1 bandit=\textsc{add} 1-be-\textsc{temp} 
that-1 alien earth.\textsc{obl-sup=add} 3-be-\textsc{temp} one trouble 3-become-\textsc{\textbf{appr}} self.\textsc{pl-dat} \textbf{whatchamacallit-\textsc{appr}} \textbf{be.afraid-\textsc{ipfv-cvb}} 
self.\textsc{pl-pl.nom} near {\textlangle}\textsc{hpl}{\textrangle}become:\textsc{pfv-neg=quot} \\
\glt `… we were afraid, there being this bandit, and us being on this land that is not ours, we were afraid that something bad may happen to us, whatchamacallit bad happens -- and did not go near there.'
\z


\noindent In \REF{ex:DD27}, the finite clause describes the action (not coming near some place) that was supposed to divert the undesirable situation (getting into trouble), but while these two situations apparently could have been described by only two clauses, as in e.g.\  \REF{ex:DD19}, here the link between the two is provided by the converb of the verb `fear', and it is unclear whether the undesirable situation instilled true emotion or just needed to be avoided.

The following example comes close to an `in-case' context; its interpretation depends on whether the speaker implies that settling down in small groups helps to avoid attracting outlaws (negative purpose interpretation) or that, while outlaws cannot be fully diverted, this reduces the damage from their attacks (`in-case' interpretation):

\begin{exe} \ex \label{ex:DD28} \textnormal{naturalistic}\\
\gll \textbf{a{\textlangle}b{\textrangle}\textsc{l}ə-lkːut} \textbf{\textsc{l}'inč'a-r-ši} jemmet bibiʁuda q'ˁe-b-di-li, jemmet t'inn-ən-n-ib bo\textsc{l} e-b-di-li\\
{\textlangle}\textbf{\textsc{hpl}{\textrangle}come-\textsc{appr}} \textbf{be.afraid-\textsc{ipfv-cvb}} thus here.and.there sit:\textsc{pfv-hpl-pst-evid} thus little-\textsc{univ-attr-pl} people be-\textsc{hpl-pst-evid}\\
\glt `Afraid that (the outlaws) would come, they were thus settling at different places, so that there were few people here and there (i.e.\ they were settling down in small groups not to attract attacks from gangs).'
\end{exe}

 Another naturalistic context relevant for the discussion is \REF{ex:DD29}:

\ea \label{ex:DD29} \textnormal{naturalistic} \\
\gll qʷat'i-li-s jatːik e{\textlangle}r{\textrangle}χ-d-inč'išː=u e{\textlangle}r{\textrangle}ku-na \textbf{he\textsc{l}'əna} \textbf{ke-lkːut} \textbf{\textsc{l}'inč'a-r} bo-li ik'-mi-s eku \\
tree-\textsc{obl-dat} up {\textlangle}2{\textrangle}climb-\textsc{pfv-cond=add} {\textlangle}2{\textrangle}fall:\textsc{pfv-temp} \textbf{thing} \textbf{4.become-\textsc{appr}} \textbf{be.afraid-\textsc{ipfv}} say:\textsc{pfv-cvb} heart-\textsc{obl-dat} 4.fall:\textsc{pfv} \\
\glt `It occurred to me: in case I climb up the tree and fall, what if something (bad) happens (to me).' (The speaker was pregnant at the moment described in the narrative, so using `something bad happens' may be viewed as a figure of avoidance).
\z

\noindent Unlike in examples \REF{ex:DD27} and \REF{ex:DD28}, in \REF{ex:DD29} the verb `fear' is in its finite form. Yet, the whole sentence is presented as a speech reporting construction which, in Archi, can extend to internal speech (`I thought that P'), and it seems unlikely that the speaker quotes her thought as saying to herself -- \textit{I thought to myself: ``You know what: I am afraid I might fall and bad things might happen''}. Though not totally impossible, it is highly improbable that she was explaining to herself her own feelings and frights. Indeed, the example has been translated by our native speaker consultant simply as independent apprehensive. In this case, the verb `fear' may be interpreted as a semantically void matrix predicate used simply to host the subordinate apprehensive clause. As a result, the whole utterance introduces the undesirable situation without introducing the way to avoid it as the main clause (though this way is indirectly alluded to by the cause that may lead to this effect shown in the conditional clause `if I climb this tree' -- in other words, don't climb the tree, and bad things won't happen to you). Semantic bleaching and ultimately omission of the verb of `fear' is thus a possible account of how the form in \nobreakdash-\textit{lkːut} underwent insubordination and developed uses outside the scope of `be afraid'. 

To sum up, we see that the use of the apprehensive in subordinate constructions covers a range of meanings, from `fear' clauses where the apprehensive shows the reason of the emotion of fear to negative purpose clauses and `in-case' clauses introducing, in the main clause, a way to avoid a negative situation or to diminish its negative impact. What all these subordinate constructions share is that they introduce a reason for the situation described in the main clause. Constructions involving the verb  `fear' may cover this whole functional range, and one naturalistic context suggests that a `fear' clause may be interpreted in a way similar to independent apprehensives.
 

\subsection{Syntactic status}
\label{subsec:Asyntactic}
\largerpage

Above we have noticed that the syntactic status of the apprehensive as the main or dependent form is sometimes ambiguous. It is relatively uncontroversial that it is independent when used in isolation, as in \REF{ex:DD1}, \REF{ex:DD2}, \REF{ex:DD3} and \REF{ex:DD5}. It is also relatively uncontroversially subordinate when used with the verb of `fear', as in \REF{ex:DD21} to \REF{ex:DD29}. But while in `fear' clauses the apprehensive clause may be analyzed as the complement clause of the verb `fear', this analysis does not work for other examples, where there is no complement taking matrix predicate and the apprehensive clause is, at best, an adverbial clause whose status is similar to purpose or reason clauses, but could sometimes just as well be viewed as a paratactic use of the independent apprehensive (cf.\ \REF{ex:DD30}). The position of the apprehensive clause with respect to the main clause does not provide a reliable clue. In Archi, a position of adverbial clauses preceding the main clause may be taken as the default, but postposition is also relatively widespread. In the case of the apprehensive, both orders can be elicited:\footnote{One of the reviewers suggested that the use of the apprehensive with the imperative (cf.\ \REF{ex:DD9}--\REF{ex:DD11} above) is also uncontroversially subordinate, because it has an explanatory function (her example is English \textit{Close the window, lest we catch a cold!}). Another reviewer, on the other hand, suggested that the use of the apprehensive with `fear' does not necessarily mean subordination. This indicates room for divergent intuitions, often based on crosslinguistic rather than language-specific evidence. We prefer to be cautious and to avoid purely functional semantic and comparative arguments in deciding upon the status of such constructions. Contra the first objection, clauses combined in a paratactical construction can easily be in explanatory relation one to another (cf.\ \textit{Close the window, we might catch a cold!}). Classifying English \textit{lest}-constructions with an imperative main clause as subordinate is based on the fact that the \textit{lest}-construction cannot be used independently, which is not true of the Archi apprehensive. As to the second objection, at least some naturalistic examples bury the apprehensional clause together with the verb of `fear' deep inside subordinate syntax. In \REF{ex:DD27}, for example, it occurs between converbial clauses and the main predicate to which they are related, which shows that the apprehensive \textit{may} be a complement of the verb of `fear'. Of course, all this evidence does not prove that it always \textit{is} its complement when they come together, but this is a general issue with syntactic evidence. What we do in this section is essentially admit that we have no syntactic evidence to decide what is the syntactic status of the apprehensive in many multi-clause constructions, including when combined with the imperative.}


\ea \label{ex:DD30} \textnormal{elicited} \\
\gll \textbf{w-eq'ːi-lkːut,} gʷači   e{\textlangle}b{\textrangle}t'ni \\
\textbf{1-bite-\textsc{appr}} dog {\textlangle}3{\textrangle}tie:\textsc{pfv} \\
\glt `They tied the dog lest it bites (him).'
\z


\ea \label{ex:DD31} \textnormal{elicited} \\
\gll gʷači-li-ɬːu iχtilat b-uq'ːi \textbf{w-eq'ːi-lkːut} \\
dog-\textsc{obl-comit} joke 3-leave.\textsc{imp} \textbf{1-bite-\textsc{appr}} \\
\glt `Do not tease the dog so that it doesn't bite you.'
\z


On the basis of the limited corpus data that we have, it seems that, as compared to the uncontroversially subordinate clauses such as those with special converbs, the apprehensive clause is more likely to be postposed to the main clause. But even this cannot be taken as a token of its main clause status, because this also happens when it is introduced by the converb form of `fear', in which case we analyze it as dependent. On the other hand, it can also be embedded into the main clause, which we also take as an unambiguous indication that the clause is subordinate (only elicited examples are available):

\ea \label{ex:DD32} \textnormal{\citet[292]{Kibrik1977}, elicited} \\
\gll to-w-mu-s zari, \textup{[}χːʷalli \textbf{bu-šbu-lkːut} deχˁʷ\textup{]} \textsc{l}o-qi \\
this-1-\textsc{obl.1-dat} I.\textsc{erg} bread 3-take-\textbf{\textsc{appr}} loaf 4.give:\textsc{pfv-fut} \\
\glt `I will give him a loaf, so that he does not buy bread.' \\ 
\z

There are also two morphosyntactic properties that group the apprehensive with non-finite forms. First, it does not have a synthetic negative counterpart. While most finite forms have a morphological way to derive negative polarity forms, most non-finite forms do not. Second, the form remains unchanged when used in uncontroversially subordinate clauses. The most common way to form subordinate constructions in Archi is the use of converbs, and none of the finite forms can be used in the context of uncontroversial subordination. We conclude that, while the syntactic status of the apprehensive clause as main vs.\ subordinate may vary across contexts, and remains unclear for many examples quoted above, morphosyntactic evidence suggests that the apprehensive groups with non-finite forms. 



\subsection{Archi apprehensive: A summary}
\label{subsec:Aapprehensive}
 
In Archi, the apprehensive is used in independent and subordinate clauses. In all of its uses, it conveys the meaning of anticipating a negative situation, including not only a warning, but also a negative purpose, events whose negative effects one has to avoid or diminish, and complements of `fear'. It is not limited to -- though often associated with -- the expression of the emotion of fear; and the lexical verb `fear' seems to be optionally available in constructions that do not imply an emotion, in contexts similar to the independent apprehensive. The exact grammaticalization source of the apprehensive remains unclear, but these are indications that its uses might have syntactically evolved from a dependent to a main clause through insubordination of a clause subordinate to a predicate of `fear' \REF{ex:DD33a}, acquiring some properties of a directive \REF{ex:DD33b}. Interestingly, in some cases the same process of insubordination is masked by the fact that the whole structure in which the main clause containing the verb of `fear' is deleted, is itself subordinate to the main clause describing the action undertaken to avoid or diminish the negative impact (\ref{ex:DD33c}):%\footnote{Note that constructions shown in \REF{ex:DD33c} are not a stage of evolution following \REF{ex:DD33b} but an independent result of dropping `fear' in \REF{ex:DD33a} when it itself is used in subordination}
 
    \ea \label{ex:DD33} \textnormal{Evolution of the apprehensive in Archi} \\
\ea \label{ex:DD33a} \textnormal{source construction} \\
\textnormal{{\ob\ob}negative effect{\cb} fear{\cb} = `fear' complement}
\ex \label{ex:DD33b} \textnormal{evolution of the independent apprehensive} \\
\textnormal{{\ob\ob}negative effect{\cb} \sout{fear}{\cb} $\rightarrow$ warning, directive} \\
\ex \label{ex:DD33c}\textnormal{evolution of the subordinate apprehensive }\\
\textnormal{{\ob\ob\ob}negative effect{\cb} \sout{fear.\textsc{cvb}}{\cb} way-to-avoid] $\rightarrow$ negative purpose, `in-case' clauses} \\
\z
\z


\noindent The constructions shown in (\ref{ex:DD33c}) should not be understood as a stage of evolution subsequent to \REF{ex:DD33b}. They independently result from dropping `fear' in a special type of context where `fear' itself was in subordinate function: where it introduced a converbial clause syntactically intermediate between the embedded apprehensive clause and the main clause into which the `fear'-converb itself was embedded. If this is true, in such contexts insubordination of the apprehensive (understood as the loss of its main predicate) did not lead to its independent use. This is what we label a \textit{masking} effect; the same process that leads to independent uses through insubordination in \REF{ex:DD33b} leads to innovative subordinate constructions in (\ref{ex:DD33c}).  

While the origins of the form are not well understood, we suggest that it was originally a subordinate form typical of the verb `fear' that evolved to an apprehensive through insubordination and loss of its matrix verb. 
 
 
 
 
 
 

\section{Mehweb: From independent apprehensive to subordinate uses via speech reporting}
\label{sec:Mehweb}
 
As in Archi, the origin of the dedicated apprehensive suffix \textit{-ala} in Mehweb is unclear, which in itself is noteworthy: Mehweb is otherwise morphologically close to other northern Dargwa varieties, where the form seems not to be attested. The first vowel of the apprehensive affix can be interpreted as the vowel of the \textit{-a-} stem also present in several forms with irrealis meaning, including optative, irrealis, and hypothetical conditional converb (see \citealt{Daniel2019,Dobrushina2019}). This would place the apprehensive in the domain of irrealis. Another possibility is a nominal origin, with a non-productive deverbal nominalization suffix \nobreakdash-\textit{ala} attested in Standard Dargwa which is based on another northern Dargwa variety. 

As in Archi, in Mehweb the apprehensive can be used in subordinate clauses. Unlike in Archi, it can be argued to be an essentially finite form that acquires subordinate uses via extension of a speech reporting construction. The corpus of Mehweb is much smaller than the corpus of Archi, and is, similarly to the corpus of Archi, exclusively narrative. However small the corpus may be, the complete absence of naturalistic examples may be telling. We will briefly come back to this point in \sectref{daniel:summary}. All discussion below is based exclusively on elicitation.


\subsection{Semantics}
\label{subsec:Msem}
 
When used in an independent clause, the apprehensive expresses the speaker's concern that an undesirable situation may come true. Imperative clauses that go together with the apprehensive may describe an action necessary to prevent it:
 

\ea \label{ex:DD34}
\gll d-arʔ-a mura, zab d-aˤq'-\textbf{ala} \\
\textsc{npl}-gather:\textsc{pfv-imp.tr} hay rain \textsc{npl}-go:\textsc{pfv-\textbf{appr}} \\
\glt `Collect the hay, it might rain.'
\z

 Unlike in Archi, the apprehensive in Mehweb shows regular negative morphology \REF{ex:DD35}. However, this cannot be used as an argument in favor of its syntactically finite rather than subordinate nature. In Dargwa languages, all forms except commands share one morphological strategy to express polarity.
 
\ea \label{ex:DD35}
\gll zab {ʜaˤ}-d-aˤq'-\textbf{ala}, d-aˤq-a šin agarod-le-ħe \\
rain \textsc{\textbf{neg}-npl-}go:\textsc{pfv-\textbf{appr}} \textsc{npl}-hit:\textsc{pfv-imp.tr} water vegetable.garden-\textsc{obl-in(lat)} \\
\glt `Let the water go to the garden, [because/in case] it might not rain.'
\z


 In \REF{ex:DD36}, the apprehensive expresses a warning about something that may happen to the addressee:
 
\ea \label{ex:DD36}
\gll q'ejju, w-igʷ-\textbf{ala} \textnormal{/} ar-d-ik-\textbf{ala} \\
slow \textsc{m}-burn:\textsc{pfv-\textbf{appr}} / \textsc{pv-f1}-fall:\textsc{pfv-\textbf{appr}} \\
\glt `Careful, don't you get burnt / don't you fall.'
\z


 Both first and third person human subjects are also available with the apprehensive. The second and first person apprehensives typically express apprehension about a potential situation whose subject the speaker \REF{ex:DD37a} or the addressee \REF{ex:DD37b} may unintentionally become. Third person apprehensives are more generally warnings about any kind of undesirable situation, and their subject may be the potential maleficiary \REF{ex:DD37c} or the agent behind the negative situation \REF{ex:DD38}, or a natural force (\REF{ex:DD34} and \REF{ex:DD35} above). In \REF{ex:DD34}, \REF{ex:DD35} and \REF{ex:DD38} the maleficiary is left unspecified and is probably understood as the addressee, at least by the pragmatic default.


\ea \label{ex:DD37}
\ea \label{ex:DD37a} \textnormal{first person} \\
\gll \textbf{nu} ʁanq' \textbf{uh-ala} \\
\textbf{I} drown \textbf{\textsc{m}.become:\textsc{pfv-appr}} \\
\glt `I might drown!' \\
\ex \label{ex:DD37b} \textnormal{second person} \\
\gll \textbf{ħu} ʁanq' \textbf{uh-ala} \\
\textbf{you.sg} drown \textbf{\textsc{m}.become:\textsc{pfv-appr}} \\
\glt `You might drown! Beware not to drown!' \\
\ex \label{ex:DD37c} \textnormal{third person} \\
\gll \textbf{eli} ʁanq' \textbf{uh-ala} \\
\textbf{child} drown \textbf{\textsc{m}.become:\textsc{pfv-appr}} \\
\glt `The child might drown!'
\z
\z


\ea \label{ex:DD38} \textnormal{third person} \\
\gll \textbf{žanawal-li-ni} maza \textbf{ar-b-uk-ala} \\
\textbf{wolf-\textsc{obl-erg}} sheep \textbf{away-\textsc{n}-lead:\textsc{pfv-appr}} \\
\glt `Beware, the wolf can steal the sheep.'
\z


 The negative evaluation is inherent to the apprehensive. When used in reference to situations which are usually viewed as positive, the evaluation necessarily changes its default value from positive to negative. Example \REF{ex:DD39} is felicitous only if the speaker wants a daughter more than a son, a preference which is culturally unexpected. Example \REF{ex:DD40} is only felicitous if the speaker does not want to recover from her illness. 
 
\ea \label{ex:DD39}
\gll urši \textbf{w-aq'-ala} ħu-ni d-aq'-a dursi \\
boy \textbf{\textsc{m}-do:\textsc{pfv-appr}} you.sg\textsc{-erg} \textsc{f1}-do:\textsc{pfv-imp.tr} girl \\
\glt `[I am afraid that] you give birth to a boy, [better] give birth to a girl!'
\z

\ea \label{ex:DD40}
\gll ara \textbf{d-uh-ala} \\
healthy \textbf{\textsc{f1}-become:\textsc{pfv-appr}} \\
\glt `[I am afraid that] I may recover!'
\z



\subsection{Subordinate uses}
\label{subsec:Msubordinate}
 
Apprehensives are used in subordination to verbs of `fear', in reference to both present \REF{ex:DD41} and past \REF{ex:DD42} situations. In these contexts, they are followed by the perfective converb of `say', \textit{ile}.


\ea \label{ex:DD41}
\gll nu uruχ k'u-we le-w-ra žanawal-li-ni maza ar-b-uk-\textbf{ala} \textbf{i-le} \\
I be.afraid \textsc{lv:ipfv-cvb} be-\textsc{m-ego} wolf-\textsc{obl-erg} sheep away-\textsc{n}-lead:\textsc{pfv-\textbf{appr}} \textbf{say:\textsc{pfv-cvb}} \\
\glt `I am afraid that a wolf will steal a sheep.'
\z

\ea \label{ex:DD42}
\gll baba uruχ k'u-we le-l-le χʷe {q'ac'} b-ik-\textbf{ala} \textbf{i-le} \\
grandmother be.afraid \textsc{lv:ipfv-cvb} \textsc{aux-f-pst} dog {bite} \textsc{3-lv:pfv-\textbf{appr}} \textbf{say:\textsc{pfv-cvb}} \\
\glt `Grandmother feared lest the dog bite (her).'
\z

 Taking the use of \textit{ile} as key evidence, we suggest that subordinate uses of the apprehensive are acquired through extending speech reporting. In East Caucasian languages, speech reporting has various extended uses, including introducing purpose. In \REF{ex:DD41} and \REF{ex:DD42}, it introduces the cause of fear. In other words, the source constructions in \REF{ex:DD41} and \REF{ex:DD42} are literally interpreted as \textit{I am afraid, saying: ``The wolf may steal a sheep!''}, \textit{Grandmother was afraid, saying: ``What if the dog bites (me)''.} It seems that, with verbs of `fear', one cannot use the apprehensive omitting \textit{ile}. An independent apprehensive is expected to refer to the moment of speech, so that, in sentences referring to the present, one could have expected to find at least a paratactic reading (`Grandmother is afraid, what if the dog bites her!'). Example \REF{ex:DD43} shows that this expectation is not borne out (omission of \textit{ile} is judged infelicitous both in reference to the past and to the present).
 

\ea \label{ex:DD43}
\gll *baba uruχ k'u-we le-r\textup{/}le-l-le χʷe \textbf{q'ac'} \textbf{b-ik-ala} \\
grandmother be.afraid \textsc{lv:ipfv-cvb} \textsc{aux-f}/\textsc{aux-f-pst} dog \textbf{bite} \textbf{\textsc{3-lv:pfv-appr}} \\
\glt (`Grandmother is/was afraid -- what if the dog bites (her).')
\z


In other subordinate constructions, the use of \textit{ile} is apparently optional. This is true not only of the combinations of the apprehensive with the imperative, as in \REF{ex:DD44}, which may or even must be interpreted paratactically (we have no data on whether \textit{ile} may be used here at all), but also of negative purpose clauses as in \REF{ex:DD45} and `in-case' clauses as in \REF{ex:DD46}, where it can be used but may be omitted:

\ea \label{ex:DD44}
\gll c'a-li-če ħule w-iz-e, \textbf{b-uš-ala} \\
fire-\textsc{obl-sup(lat)} look \textsc{m-lv:pfv-imp} \textbf{\textsc{n}-die(of.fire):\textsc{pfv-appr}} \\
\glt `Watch the fire, it might go out.'
\z


\ea \label{ex:DD45}
\gll it-ini ħuli b-uc-ib c'a \textbf{b-uš-ala} \textbf{\textup{(}i-le\textup{)}} \\
this-\textsc{erg} look \textsc{n-lv:ipfv-pst} fire \textbf{\textsc{n}-die.out:\textsc{pfv-appr}} \textbf{(say:\textsc{pfv-cvb})} \\
\glt `He was looking after the fire so it would not die out.'
\z


\ea \label{ex:DD46}
\gll Musa-ni mura d-arʔ-ib dunijal \textbf{ur-ala} \textbf{\textup{(}i-le\textup{)}} \\
Musa-\textsc{erg} hay \textsc{npl}-gather:\textsc{pfv-aor} world \textbf{rain:\textsc{ipfv-appr}} \textbf{(say:\textsc{pfv-cvb})} \\
\glt `Musa collected the hay out of fear that rain would start' (lit.\ `lest the world rains').
\z


 Another way to express negative purpose, alternative to the apprehensive $+$ \textit{ile} strategy in \REF{ex:DD46}, is the use of the purposive converb in \textit{-alis}. Unlike with the apprehensive, where the negative evaluation is hardwired into the semantics of the form itself, the negative component `not to' is overtly expressed on the purposive converb by the prefix \textit{ħa-} (standard negation). The converb \textit{ile} cannot be used with the purposive converb:


\ea \label{ex:DD47}
\ea \label{ex:DD47a}
\gll w-aˤld-e adaj-ni ħu dam \textbf{ħa-w-aq'-a-lis} \\
\textsc{m}-hide:\textsc{pfv-imp} father-\textsc{erg} you.sg beat \textbf{\textsc{neg-m}-do:\textsc{pfv-irr-purp}} \\
\glt `Hide, so that your father does not beat you.' \\
\ex \label{ex:DD47b}
\gll w-aˤld-e adaj-ni ħu dam \textbf{w-aq'-ala} \\
\textsc{m}-hide:\textsc{pfv-imp} father-\textsc{erg} you.sg beat \textbf{\textsc{m}-do:\textsc{pfv-appr}} \\
\glt `Hide, so that your father does not beat you.'
\z
\z


 Above, we have argued for Archi that what is common for the use of the apprehensive in all subordinate clauses -- including `fear' complement clauses, negative purpose clauses and `in-case' clauses -- is that the apprehensive provides a (negative) reason for the action described in the main clause. This becomes very clear when comparing the scope of the use of the apprehensive to that of the negative purpose converb in Mehweb. While the purpose converb competes with the apprehensive in  contexts where both describe an undesirable situation that the subject of the main clause may control, the purpose converb is not available in  clauses which describe situations out of her control. For example,  \REF{ex:DD48}, an `in-case' clause, and  \REF{ex:DD49}, a `fear' clause, are pragmatically infelicitous.

\ea \label{ex:DD48}
\gll *Musa-ni mura d-arʔ-ib dunijal \textbf{ħa-wr-a-lis} \\
Musa-\textsc{erg} hay \textsc{npl}-gather:\textsc{pfv-aor} world \textbf{\textsc{neg}-rain:\textsc{ipfv-irr-purp}} \\
\glt Intended (same as \REF{ex:DD46}):
 `Musa collected the hay so that it would not rain.'
\z


\ea \label{ex:DD49}
\gll *baba uruχ k'u-we le-l-le χʷe \textbf{q'ac'} \textbf{ħa-b-ik-a-lis} \\
grandmother be.afraid \textsc{lv:ipfv-cvb} \textsc{aux-f-pst} dog \textbf{bite} \textbf{\textsc{neg-3-lv:pfv-irr-purp}} \\
\glt Intended (same as \REF{ex:DD42}):  `Grandmother became afraid so as to make the dog not bite her.'
\z


\noindent In \REF{ex:DD48} and \REF{ex:DD49}, we see that the purpose converb is not available in some contexts where the apprehensive with \textit{ile} can be used (\REF{ex:DD46} and \REF{ex:DD42}, respectively). While the purpose reading for the apprehensive construction is only a possible contextual implication of its more general semantics of reason, the purpose converb is dedicated to expressing a purpose of an action, which neither a stimulus of fear nor a meteorological situation can (normally) be.

In addition to its uses as warning, the apprehensive thus also covers subordinate uses described above for Archi, including `fear' clauses, negative purpose clauses (where it competes with the dedicated purpose converb in its negative form) and `in-case' clauses. However, apprehensive `fear' clauses must, and apprehensive negative purpose clauses may, be followed by the converb of the verb `say', a peculiar property that will be central to the interpretation of its syntactic status in the following section.


\subsection{Syntactic status}
\label{Msyntactic}
 
The syntactic status of the apprehensive in Mehweb is not straightforward. As in Archi, the apprehensive clause in Mehweb can be used both in isolation and in polypredicative constructions, but in the latter case it is often followed by the converb of the verb of speech, \textit{ile}. When combined with an imperative, the apprehensive is apparently used without further morphosyntactic complications, as when used in isolation (we do not have data on whether \textit{ile} can even be used in such contexts in the first place). Imperative-apprehensive contexts may thus be interpreted as paratactic constructions rather than constructions involving subordination. In more clear cases of subordination, an apprehensive clause may or even must be followed by \textit{ile.} Thus, with verbs of `fear', one cannot use the apprehensive omitting \textit{ile}, while negative purpose and `in-case' clauses allow for both constructions. 

To explain the suggested evolution of the subordinate uses of the Mehweb apprehensive, it is necessary to provide a background on speech reporting in Mehweb, by and large typical of speech reporting in East Caucasian. Contexts of speech reporting are the most common use of \textit{ile}, literally `having said'. Cf.\ the following example: 
 
\ea \label{ex:DD50} \textnormal{naturalistic}\\
\gll sune-la=l hubi-ze iχ-ini ``\textbf{ma,} \textbf{ħa-la=l} \textbf{urši''} \textbf{i-le,} kʷe-w u-ub-i urši hub-la q'uq'u-be-če w-ix-ib\\
self.\textsc{obl-gen=ptcl} husband-\textsc{inter(lat)} that-\textsc{erg} \textbf{take} \textbf{you.sg.\textsc{obl-gen=ptcl}} \textbf{boy} \textbf{say:\textsc{pfv-cvb}} in.hands-\textsc{m(ess)}	\textsc{m}.become:\textsc{pfv-aor-ptcp} boy husband-\textsc{gen} knee-\textsc{pl-sup(lat)} \textsc{m}-put:\textsc{pfv-aor}\\
\glt `She said to her husband: ``Here you go, (this is) your son'' and put the boy, whom she had in her hands, onto his lap.'
\z

\noindent In \REF{ex:DD50}, the speech reporting clause is embedded between the converb of the verb of speech and its argument (`she ``\_'' having said, \ldots'). Speech reporting thus seems to form a subordinate clause, conveying an indirect report. The converb of `say' may combine with a lexical speech predicate, as in \REF{ex:DD51} and \REF{ex:DD52}, where it may seem to function as a complementizer.

\ea \label{ex:DD51} \textnormal{naturalistic} \\
\gll urši-li-ʔini \textbf{žawab} \textbf{g-ib} k'asːi-ze ``ħu-ni b-aqˁ-a'' \textbf{i-le} \\
boy-\textsc{obl-erg} \textbf{answer} \textbf{give:}\textsc{pfv-aor} fish-\textsc{inter(lat)} you.sg.\textsc{erg} \textsc{n}-hit-\textsc{imp.tr} say:\textsc{pfv-cvb} \\
\glt `The youth answered to the fish: ``You hit''!'
\z

\ea \label{ex:DD52} \textnormal{naturalistic} \\
\gll nu-ni \textbf{xar} \textbf{b-a-i-ra} ``hej urši, kuda-b-a heš ulica'', \textbf{i-le} \\
I-\textsc{erg} \textbf{request} \textsc{n-lv:pfv-aor-ego} hey boy where-\textsc{n(ess)-intrg} that street \textbf{say:\textsc{pfv-cvb}} \\
\glt `I asked: ``Hey guy, where is this street?'''
\z

 There are several arguments against such an analysis. First, in Mehweb, a language with otherwise strictly non-finite subordination, speech reporting would represent the only construction which involves a finite clause in an apparently subordinate function. Second, reported speech introduced by \textit{ile} may involve imperatives, such as `hit!' in \REF{ex:DD51} and expressions of address, such as `hey guy!' in \REF{ex:DD52}, which are all main clause phenomena, as well as preserve original deixis, e.g.\  `you' in \REF{ex:DD51} and `your' in \REF{ex:DD50}. All this can be seen as indication that the speech reporting strategy using \textit{ile} is direct rather than indirect,\footnote{Whether indirect reporting does exist in East Caucasian languages at all is in fact a controversial issue.} with the reporting quote exemplifying what \citet{SpronckNikitina2019} call a \textit{dedicated syntactic domain}, which explains both the availability of main clause phenomena and the use of finite forms. The use of \textit{ile} in speech reporting constructions, while frequent, is not obligatory, as shown in \REF{ex:DD53} (compare with \REF{ex:DD51}) and \REF{ex:DD54} (compare with \REF{ex:DD52}).

\ea \label{ex:DD53} \textnormal{naturalistic} \\
\gll bars \textbf{žawab} \textbf{g-ib:} ``w-ig-an!'' \\
in.exchange \textbf{answer} \textbf{give:\textsc{pfv-aor}} \textsc{m}-want-\textsc{prs} \\
\glt `They answered back: ``(We) want him!''{}'
\z


\ea \label{ex:DD54} \textnormal{naturalistic} \\
\gll nu-ni hanna=ra \textbf{xar} \textbf{b-a-i-ra:} ``hej urši, heš ʁʷada-ni saʁʷa uˤq'-es-i?'' \\
I-\textsc{erg} now=\textsc{add} \textbf{request} \textbf{3-\textsc{lv:pfv-aor-ego}} hey boy this street-\textsc{loc(lat)} how \textsc{m}.go-\textsc{inf-attr} \\
\glt `I ask again: ``Hey, boy, how should one go to this street?''{}'
\z


 To sum up the relevant data on speech reporting, in Mehweb, what seems a subordinate indirect speech reporting clause in \REF{ex:DD50}, is in fact embedded direct speech reporting. We suggest that the uses of \textit{ile} to introduce the apprehensive clause may be viewed as an extension of direct speech reporting constructions and do not contradict the status of the apprehensive as a finite form. 

A collateral argument comes from reference tracking. The use of the \textit{sa-CL-i} `self' pronoun in \REF{ex:DD55} is indicative of the status of the construction as extended speech reporting, because the `self' pronoun is here used in its logophoric function, under the typical reported speech condition of coreference of an argument of the dependent clause with the subject of the matrix clause.


\ea \label{ex:DD55}
\gll baba uruχ k'-uwe le-r, \textbf{sa{\textlangle}r{\textrangle}i} \textbf{ar-d-ik-ala} \textbf{i-le} \\
granny be.afraid \textsc{lv:ipfv-cvb} \textsc{cop-f} \textbf{self{\textlangle}\textsc{f.sg}{\textrangle}} \textbf{\textsc{pv-f1}-fall:\textsc{pfv-appr}} \textbf{say:\textsc{pfv-cvb}} \\
\glt `The grandmother is afraid of falling.'
\z


This supports our suggestion that apprehensive clauses with \textit{ile} are extensions of speech reporting. It also explains how some of the restrictions on the use of the apprehensive are removed. In isolation, the apprehensive is used from the temporal and evaluative perspective of the current speaker. Speech reporting introduces a new perspective, that of the reported speaker, which includes an independent frame of temporal reference and evaluative basis, including allowing for the actual speaker not to share the evaluation of the subject of the main clause. What seems to be an apprehensive referring to the past is in fact a regular apprehensive referring to the present but in a new system of temporal reference shifted to the viewpoint of the subject of the main clause. 

As yet another piece of evidence for the essentially non-subordinate nature of the apprehensive, consider the following example. In \REF{ex:DD56}, where the apprehensive clause is followed by \textit{ile}, it can be embedded.

\ea \label{ex:DD56}
\gll Musa-ni \textup{[}dunijal \textbf{ur-ala} \textbf{i-le}\textup{]} mura d-arʔ-ib \\
Musa-\textsc{erg} world \textbf{rain:\textsc{ipfv-appr}} \textbf{say:\textsc{pfv-cvb}} hay \textsc{npl}-gather:\textsc{pfv-aor} \\
\glt `Musa collected the hay out of fear that rain starts.'
\z

 Without the \textit{ile} converb, the apprehensive clause cannot be embedded. Cf.\ the following examples, ungrammatical with, but grammatical without embedding:


\ea \label{ex:DD57}
\ea[]{ \label{ex:DD57a}
\gll eli šula-le b-uc-a \textup{[}ʁadara \textbf{b-oˁrʡ-aq-ala}\textup{]}\\
child tight-\textsc{advz} \textsc{n}-hold:\textsc{pfv-imp.tr} dish \textbf{\textsc{n}-break:\textsc{pfv-caus-appr}}\\
\glt `Hold the child tight so that it does not break the dish.' 
}
\ex[*]{ \label{ex:DD57b}
\gll eli \textup{[}ʁadara \textbf{b-oˁrʡ-aq-ala}\textup{]} šula-le b-uc-a \\
child dish \textbf{\textsc{n}-break-\textsc{caus-appr}} tight-\textsc{advz} \textsc{n}-hold:\textsc{pfv-imp.tr} \\
\glt Intended: `Hold the child tight so that it does not break the dish.' 
}
\z
\z

\ea \label{ex:DD58}
\ea[]{ \label{ex:DD58a}
\gll sumka b-ux-a mataħ \textbf{ar-d-uʔ-ala} \\
bag \textsc{n}-bring:\textsc{pfv-imp.tr} money \textbf{\textsc{pv-npl}-lose:\textsc{pfv-appr}} \\
\glt `Take the bag so you don't lose the money.'\\
}
\ex[*]{ \label{ex:DD58b}
\gll sumka \textup{[}mataħ \textbf{ar-d-uʔ-ala}\textup{]} b-ux-a \\
bag money \textbf{\textsc{pv-npl-}lose:\textsc{pfv-appr}} \textsc{n}-bring:\textsc{pfv-imp.tr}\\
\glt Intended: `Take the bag so you don't lose the money.'
}
\z
\z

 This behavior is parallel to speech reporting. When a reporting quote follows the speech clause, it may (\REF{ex:DD51}, \REF{ex:DD52}) and may not (\REF{ex:DD53}, \REF{ex:DD54}) use \textit{ile.} When it is embedded, as in \REF{ex:DD50}, dropping \textit{ile} -- which also serves as a delimiter for Spronck and Nikitina's dedicated syntactic domain of reporting -- is supposedly less frequent (unfortunately, we do not have the relevant data to argue it is ungrammatical).  

We conclude that without \textit{ile} the form is an independent apprehensive, and in all such contexts where it is used in a polypredicative environment (probably limited to the use with imperatives, as in \REF{ex:DD57a} and \REF{ex:DD58a}) it should be interpreted as paratactic. In the contexts that are more likely to be subordinate, as `in-case' clauses, negative purpose clauses and `fear' clauses, its use evolves as a functional extension by adopting the syntactic pattern of direct speech reporting. 



\subsection{Mehweb apprehensive: A summary}
\label{subsec:Mapprehensive}
 
As in Archi, Mehweb apprehensives have both independent and subordinate uses. In subordinate uses they cover functions similar to those in Archi, including negative purpose, `in-case' clauses and complements of `fear'. Unlike in Archi, subordinate uses seem to be associated with the use of the converb \textit{ile} `having said', optionally or obligatorily, also used to introduce reported speech with verbs of speech. We suggest that the Mehweb apprehensive originates from a form in the main clause with a directive/emotive illocution that has acquired subordinate uses by functional expansion of speech reporting. The evolution of the uses of the apprehensive in Mehweb may then be summarized as follows (note that, unlike other types of subordination, speech reporting typically places quotes after the matrix verb):


\ea \label{ex:DD59} \textnormal{Evolution of the apprehensive in Mehweb} \\
\ea \label{ex:DD59a}\textnormal{source construction} \\
\textnormal{{\ob}apprehensive{\cb} = warning/directive}
\ex \textnormal{evolution of subordinate uses via extended speech reporting:} \\
\textnormal{{\ob}finite verb {\ob}apprehensive{\cb} \textit{ile}{\cb}  = `fear', negative purpose and `in-case' clauses}
\ex \textnormal{specification of some of the subordinate uses:} \\
\textnormal{{\ob}finite verb {\ob}apprehensive{\cb} \sout{\textit{ile}}{\cb}  = negative purpose and `in-case' clauses}
\z
\z

Although the origin of the apprehensive in Archi is unclear, there are some indications that it has followed the opposite path from a dependent to main clause through insubordination, acquiring some properties of a directive (cf.\ the scheme in \REF{ex:DD33} above).








\section{Northern Akhvakh: From a particle to a morphological apprehensive}
\label{sec:NAkhvakh}
 
To our knowledge, the existence of an apprehensive in Akhvakh was first reported by Denis Creissels:\footnote{While this study quotes the forms as \textit{-alogo}, our own data and later work by Denis Creissels \citeyear{Creisselsunpublished2} has \textit{-alago}.} 

\begin{quote}
    I have also observed the sporadic occurrence of \textit{an emphatic form of the prohibitive} expressing a threat in case the prohibition is not respected. This form, for which my main informant gives the Russian equivalent `Не дай Бог …!' [`God forbid' -- MD \& ND], is marked by a suffix \textit{-alogo} (\ldots). \citep{Creisselsunpublished1}
\end{quote}

Creissels collected his data in the village of Axaxdərə near Zaqatala (Azerbaijan), a resettlement from Daghestanian villages speaking Northern Akhvakh. There is no reliable information about the time when the Akhvakh started migrating to Azerbaijan. The oral tradition suggests that the migration was progressive rather than punctual, and started about one century ago (p.c.\  with Denis Creissels). The last wave of migration occurred immediately after World War 2. According to Creissels, Axaxdərə Akhvakh is still very close to the dialect spoken in the Akhvakh district in Daghestan. 

In 2018, one of the authors of this paper worked in Tad-Magitl, a Northern Akhvakh village in the Akhvakh district, and later with Indira Abdulaeva, a native of Tad-Magitl. Most examples in this section come from these field notes.

\citet{Creisselsunpublished2} suggests that the apprehensive suffix \textit{-alago} originated from the conditional converb in \nobreakdash-\textit{ala} in combination with the interjection \textit{hologo \textasciitilde hʷelogo} `beware!' . The etymology of the interjection itself is unknown. Creissels provides examples of the apprehensive and the interjection used together (for the sake of consistency, we gloss his ``prohibitive'' as apprehensive):

\ea \label{ex:DD60}
\gll hʷelogo imaχa-ge \textsc{l}'a \textbf{ʁad{\textlangle}u{\textrangle}k'-alago!} \\
beware donkey-\textsc{loc} on \textbf{\textsc{{\textlangle}m{\textrangle}}sit-\textsc{appr}} \\
\glt `Beware, don't sit on the donkey!' 
\citep{Creisselsunpublished1}
\z

\ea \label{ex:DD61}
\gll hʷelogo ha-be \textbf{gʷ-eːlago}! \\
beware \textsc{dem-n} \textbf{do-\textsc{appr}} \\
\glt `Beware, don't do this!' \citep{Creisselsunpublished1}
\z

 Our data (in the case of Akhvakh, elicited except where indicated otherwise) confirm these examples for Tad-Magitl, a Northern Akhvakh village in Daghestan. Our consultant, Indira Abdulaeva, accepts the use of the affix both with (\REF{ex:DD62}, \REF{ex:DD63}) and without \REF{ex:DD64} the particle, and notices that the apprehensive may have realizations ``extended'' by \nobreakdash-\textit{le}, as in \REF{ex:DD63} and \REF{ex:DD64}, and that the particle may occur both as \textit{hole} \REF{ex:DD62} and as \textit{holego} \REF{ex:DD63}.

\ea \label{ex:DD62}
\gll hole χʷe-de \textbf{q'ːeleč'-alago} \\
beware dog-\textsc{erg} \textbf{bite-\textsc{appr}} \\
\glt `Beware, the dog might bite you.'
\z

\ea \label{ex:DD63}
\gll Aħmadi holego \textbf{aq'ːusː-alagole!} \\
Ahmad beware \textbf{drown-\textsc{appr}} \\
\glt `Ahmad, beware, you may drown!'
\z

\ea \label{ex:DD64}
\gll c'ːa \textbf{c'ː-alagole\textup{/}c'ː-alago} \\
rain \textbf{rain-\textsc{appr}} \\
\glt `It might rain.'
\z

 In addition to \textit{hole(go)}, the dictionary of Northern Akhvakh (\citealt{MagomedovaAbdulaeva2007}) quotes another particle associated with apprehension, \textit{čego}, translated as `what if'.\footnote{Reviewers note that, in this particle, the segment \textit{go} appears once again, but even with a similar variation between \textit{hole} and \textit{holego} we have no evidence for its meaningful morphological segmentation. As a possible parallel, the dedicated apprehensive particle \textit{vore} in Avar, which is the local lingua franca, also appears in the variant \textit{voregi}, where \nobreakdash-\textit{gi} may be identified as the additive particle `also', `too', that has a wide range of discourse uses. No such evidence for \nobreakdash-\textit{go} in Akhvakh is currently available.} In its three occurrences in the dictionary, the particle is combined with the apprehensive, as in \REF{ex:DD65}; \REF{ex:DD66} is a similar example obtained in elicitation.


\ea \label{ex:DD65} \textnormal{dictionary example} \\
\gll čego mik'i-ga q'ːʷat'ala \textbf{b-eq'-alagole} \\
beware child-\textsc{loc.all} cold \textbf{\textsc{n}-come-\textsc{appr}} \\
\glt `What if the child catches a cold!'
\z

\ea \label{ex:DD66}
\gll eqː-a čego χʷe-de mik'i-ge \textbf{q'ːeleč'-alagole} \\
look-\textsc{imp} beware dog-\textsc{erg} child-\textsc{loc} \textbf{bite}-\textsc{appr} \\
\glt `Beware lest the dog bites the child!'
\z





\subsection{Semantics and morphology}
\label{subsec:NAsem}
 
The form in \nobreakdash-\textit{alago}(\textit{le}) can be used with subjects of all persons. It is thus more appropriately described as apprehensive than as emphatic prohibitive.

\ea \label{ex:DD67}
\ea \label{ex:DD67a} \textnormal{1st person} \\
\gll dene \textbf{aq'ːusː-alago\textup{(}le\textup{)}} \\
I \textbf{drown-\textsc{appr}} \\
\glt `What if I drown?' \\
\ex \label{ex:DD67b} \textnormal{2nd person} \\
\gll mene \textbf{aq'ːusː-alago\textup{(}le\textup{)}} \\
you.sg \textbf{drown-\textsc{appr}} \\
\glt `What if you drown!' \\
\ex \label{ex:DD67c} \textnormal{3rd person} \\
\gll hudu-we \textbf{aq'ːusː-alago\textup{(}le\textup{)}} \\
that-\textsc{m} \textbf{drown-\textsc{appr}} \\
\glt `What if he drowns!'
\z
\z

 The apprehensive clause may describe either the undesirable situation itself, as in \REF{ex:DD68}, or a situation that implies further negative consequences, as in \REF{ex:DD69}; probably, the two uses cannot be consistently separated, because a situation that implies negative outcomes is itself perceived as negative. 


\ea \label{ex:DD68}
\gll ɬːeni goɬiːhi ba-χːʷ-alagole i\textsc{l}ːi \\
water be:\textsc{neg:hpl:cvb} \textbf{\textsc{hpl}-remain-\textsc{appr}} \textsc{incl} \\
\glt `We may be left with no more water!'
\z


\ea \label{ex:DD69}
\gll \textsc{l}'ːõhe \textbf{W-ux-alagole} mene \\
fall.asleep:\textsc{m:cvb} \textbf{\textsc{m}-become-\textsc{appr}} you.sg \\
\glt `Beware of falling asleep!'
\z


 Like in Archi, but unlike in Mehweb, the apprehensive does not have a (morphological) negative counterpart. We have no data as to how apprehension that a positively evaluated situation may not come true is expressed in Northern Akhvakh.
 

\subsection{Syntactic status}
\label{subsec:NAsyntactic}
 
According to our field data, the form in \nobreakdash-\textit{alago}(\textit{le}) introduces an independent clause and can not be used in subordination. 

In four out of the six total occurrences in \citet{MagomedovaAbdulaeva2007}, the apprehensive clause occurs in isolation. In the two other examples, \REF{ex:DD70} and \REF{ex:DD71}, it occurs in combination with an imperative; cf.\ also \REF{ex:DD72} obtained in elicitation. 

\ea \label{ex:DD70} \textnormal{dictionary example} \\
\gll imiχi \textbf{aq'ːusː-alagole} k'oli-g-une  kʷari  b-eqː-a  \\
donkey \textbf{drown-\textsc{appr}} neck-\textsc{loc-el} rope \textsc{n}-take.away-\textsc{imp} \\
\glt `Take the rope off the donkey's neck, it may suffocate.'
\z

\ea \label{ex:DD71} \textnormal{dictionary example} \\
\gll χːeχːi-meč'i\textsc{l}a, \textbf{w-ač'aq'-alagole} \\
hurry-hurry \textbf{\textsc{m}-be.late-\textsc{appr}} \\
\glt `Hurry, or you'll be late!'
\z

\ea \label{ex:DD72}
\gll čadi b-eχ-a ã\textsc{l}'o \textbf{c'ː-alagole} \\
umbrella \textsc{n}-carry-\textsc{imp} together \textbf{rain-\textsc{appr}} \\
\glt `Take an umbrella, it might rain!'
\z  


\noindent Example \REF{ex:DD70} looks as if the apprehensive could be interpretable as a negative purpose clause. However, this example, as well as \REF{ex:DD71} and \REF{ex:DD72}, all contain imperatives. If the apprehensive were a negative purpose clause subordinate to `take' in \REF{ex:DD70}, it could also be used with the main verb in the past tense. This is not the case; \REF{ex:DD73}, which is a modified version of \REF{ex:DD70} with the presumably main verb in the past, is ruled out. 

\ea \label{ex:DD73}
\gll *imiχi \textbf{aq'ːusː-alagole} k'oli-g-une kʷari b-eqː-ari \\
donkey \textbf{drown-\textsc{appr}} neck-\textsc{loc-el} rope \textsc{n}-take.away-\textsc{pst} \\
\glt Intended: `(He) took the rope off the donkey's neck, lest it would get strangled.' 
\z


 Our consultant notices that the form conveys apprehension about a future event, and one cannot warn someone about the past. Clauses introducing negative purpose in the past are based on the negative converb \REF{ex:DD74a}; this construction is also available for the main clause with reference to the moment of speech, with the main verb in the present \REF{ex:DD74b}:



\ea \label{ex:DD74}
\ea \label{ex:DD74a}
\gll baba-de \textbf{aq'ːusː-i\textsc{l}-uk'u} waša ut-i\textsc{l}a \~ihʷari-ga \\
mother-\textsc{erg} \textbf{drown-\textsc{neg-purp}} son \textsc{m}.let-\textsc{pfv.neg} lake-\textsc{loc.lat} \\
\glt `The mother did not let the child go to the lake so that he would not drown.' \\
\ex \label{ex:DD74b}
\gll baba-de \textbf{aq'ːusː-i\textsc{l}-uk'u} waša ut-ero gu\textsc{l}a \~ihʷari-ga \\
mother-\textsc{erg} \textbf{drown-\textsc{neg-purp}} son \textsc{m}.let-\textsc{m.ipfv} \textsc{cop.neg} lake-\textsc{loc.lat} \\
\glt `The mother does not let the child go to the lake so that he does not drown.'
\z
\z


 The apprehensive is equally impossible in `fear' complement clauses,\footnote{One of the reviewers asks whether the apprehensive can or cannot be introduced by a quotative particle in a type of extended speech reporting construction. This would be an important piece of evidence which we very unfortunately lack.} where a locative form of the action nominal is used instead:


\ea \label{ex:DD75}
\gll dene \textsc{l}oːhe gudi du-ge χʷe-de \textbf{q'ːeleč'a-ro-g-une} \\
I be.afraid.\textsc{cvb} \textsc{cop.m} you.sg\textsc{-loc} dog-\textsc{erg} \textbf{bite-\textsc{nmlz.obl-loc-el}} \\
\glt `I'm afraid that the dog might bite you.'
\z


As for constructed contexts with a `fear' verb and an apprehensive, our consultant suggested that what results in \REF{ex:DD76} are two separate sentences, which suggests a paratactic structure, similarly to the combination of the apprehensive with an imperative in \REF{ex:DD70}, \REF{ex:DD71} and \REF{ex:DD72} above:


\ea \label{ex:DD76}
\gll dene \textsc{l}oho gudi, du-ge χʷe-de \textbf{q'ːeleč'-alago} \\
I afraid \textsc{cop.m} you.sg\textsc{-loc} dog-\textsc{erg} \textbf{bite-\textsc{appr}} \\
\glt `I'm frightened! What if the dog bites you!'
\z



 It remains unclear how \REF{ex:DD77} should be interpreted, with the apprehensive clause sitting between the conditional clause and the imperative:


\ea \label{ex:DD77}
\gll eχːa w-ano w-\~uč-ala \textbf{c'ː-alagole} čadi b-eχ-a \~an\textsc{l}'o \\
outside \textsc{m}-go.\textsc{cvb} \textsc{m}-find-\textsc{cond} \textbf{rain-\textsc{appr}} umbrella \textsc{n}-carry-\textsc{imp} together \\
\glt `If you plan to go outside, it might rain! Get an umbrella.'
\z



\noindent Presumably, the condition `if you go out' is subordinate to the imperative `get an umbrella', and it is unclear how the apparently paratactic apprehensive makes it to the position between them. More data is needed here.
 

\subsection{Akhvakh apprehensive: A summary}
\label{subsec:NAapprehensive}
 
We conclude that, in Akvakh, the apprehensive is functionally similar to the independent uses of the apprehensive in Archi and Mehweb, including availability to subjects of all persons. At the same time, it is very different from the apprehensives in both languages in terms of its syntactic behavior, because it is not used in subordination. The typical contexts of use of the apprehensive in subordinate clauses in Archi and Mehweb, including negative purpose or `fear' complement clauses, are covered by other forms, and what appears to be subordinate uses of the apprehensive has been argued to be paratactical, similar to the use of the apprehensive with imperatives in Mehweb. This is in accordance with Creissels' (\citeyear{Creisselsunpublished2}) suggestion that the apprehensive affix originated from univerbation of the verb with the interjection \textit{hologo} `beware!', as interjections are known to belong to main-clause phenomena. 


\ea \label{ex:DD78} \textnormal{Evolution of the apprehensive in Akhvakh}
\ea \label{ex:DD78a} \textnormal{source construction:} \\
\textnormal{{\ob}verb apprehensive.particle{\cb} = warning/directive}
\ex \label{ex:DD78b} \textnormal{univerbation:} \\
\textnormal{{\ob}verb-\textsc{appr}{\cb} = warning/directive}
\z
\z

 However, the hypothesis of the evolution of the apprehensive through univerbation with the particle needs to be carefully assessed. It is unclear why the particle would follow the verb (in all examples, the ``free'' particle is used sentence-initially). It is unclear why it should attach to the conditional converb. It is also unclear where the ``extended'' version \textit{-alagole} comes from, and what its etymological relation to \textit{-alago} is. Answering all these questions requires re-evaluation of the data, which may lead to reconsidering the initial hypothesis.



\section{Summary discussion}
\label{daniel:summary}
 
The apprehensive forms described above for three languages of Daghestan, Archi, Mehweb and Akhvakh, all belonging to different branches of the family, were shown to be diverse in terms of their internal structure and possible origins, but, to some extent, convergent in their functional syntactic scope. We summarize the convergent and divergent properties of the three forms in \tabref{tab:summary}.
 


\begin{table}[h]

\caption{Properties of Archi, Mehweb and Akhvakh apprehensives}
\label{tab:summary}
\begin{tabularx}{\textwidth}{QccC}
\lsptoprule
\textbf{} & \textbf{Archi} & \textbf{Mehweb} & \textbf{Akhvakh} \\
\midrule
morphological source & ? & ? & ?\newline (conditional + interjection?) \\
\tablevspace
main clause:\newline apprehensive proper & yes & yes & yes \\
\tablevspace
subordinate:\newline `fear' complement & yes & yes & no \\
\tablevspace
subordinate:\newline negative purpose & yes & yes & no \\
\tablevspace
subordinate:\newline `in-case' clauses & yes & yes & no \\
\tablevspace
person reference & 1/2/3 & 1/2/3 & 1/2/3 \\
\tablevspace
finiteness & non-finite & finite & finite \\
\tablevspace
negative form & periphrastic & morphological & unavailable \\
\lspbottomrule
\end{tabularx}
\end{table}

 
\noindent On the formal side, the apprehensive has a negative form in Mehweb, is negated by means of a periphrastic construction in Archi, and has no negative counterpart in Akhvakh at all. The Akhvakh apprehensive probably goes back to the combination of a verbal form with an apprehensive interjection still present in the languages, while the origins of the apprehensives in Mehweb and Archi are unclear. 

All of these forms express apprehension when used in the main clause, and apply to all subjects. But unlike Akhvakh, both Archi and Mehweb forms are used in subordination, including `fear' complements, negative purpose clauses, and what \citet{lichtenberk1995} has termed `in-case' clauses. Sentential `fear' complements provide the stimulus for fear, which may be viewed as an instantiation of the concept of cause or reason. We have suggested that what negative purpose (`so as not to') and `in-case' (`out of fear of') clauses have in common is that the main clause suggests a way to avoid or reduce the negative impact of the event described in the apprehensive clause, so that, in these contexts too, apprehensive constructions may be seen as reason rather than purpose adverbial clauses. All apprehensive-marked subordinate clauses may thus be interpreted as expressing the reason for another event/action (expressed by the main clause) that is carried out to avoid an undesirable consequence. But with all these apparent similarities, Archi and Mehweb are different in terms of the syntactic patterns involved. 

In East Caucasian languages, there is a strong morphological divide between main clauses using finite forms and dependent clauses using various types of converbs, infinitives and action nominals. Some forms considered primarily non-finite can be used as main clause predicates, but this happens under special circumstances. Thus, the perfective converb is used finitely for unwitnessed  events \citep{Kibrik1977}, and participles are sometimes used finitely, a phenomenon that is clearly understudied (\citealt{KalininaSumbatova2007,Kalinina2011}). In Mehweb, too, converbs can be used in main clauses with subtle pragmatic effects \citep{Kustova2019}. Conversely, finite forms may be used in apparently dependent contexts, but this happens in a special environment of extended speech reporting, such as in purpose clauses accompanied by the converb of the verb of speech, \textit{boli} in Archi and \textit{ile} in Mehweb.  

In Mehweb, the use of the apprehensive in `fear' complement clause requires the use of \textit{ile} `having said', and its use in negative purpose and `in-case' clauses allows (but does not require) it. Archi has a very similar pattern of extended speech reporting, but the apprehensive shows no association with the verb of speech -- instead, in both negative purpose and `in-case' clauses the verb `fear' can be used. In Archi, the apprehensive, unlike the finite forms but like some (probably most) non-finite forms, does not have a morphological negative counterpart. The negative correlate of the apprehensive (`lest P does \textit{not} happen') is built periphrastically, using a negative converb of the main verb with the apprehensive of the verb `remain'. 

We hypothesize that the Archi apprehensive is essentially a non-finite form that has acquired main clause uses via insubordination \citep{evans2007insubordination}, by losing its main `fear' clause. Subordinate uses may also have followed this path, but here insubordination is masked by the fact that this insubordination itself happened within a subordinate clause. In other words, the verb `fear' was dropped as both main and subordinate matrix clause dominating the apprehensive clause, leading to different outcomes:

\ea \label{ex:DD79} \textnormal{Evolution of the apprehensive in Archi (repeated from \REF{ex:DD33} above)} \\
\ea \label{ex:DD79a} \textnormal{source construction} \\
\textnormal{{\ob\ob}negative effect{\cb} fear{\cb} = `fear' complement}
\ex \label{ex:DD79b} \textnormal{evolution of the independent apprehensive} \\
\textnormal{{\ob\ob}negative effect{\cb} \sout{fear}{\cb} $\rightarrow$ warning, directive}  \\
\ex \label{ex:DD79c} \textnormal{evolution of the subordinate apprehensive} \\
\textnormal{{\ob\ob\ob}negative effect{\cb} \sout{fear.cvb}{\cb} way-to-avoid] $\rightarrow$ negative purpose, `in-case' clauses} \\
\z
\z

 In Mehweb, on the other hand, the form is probably originally a finite form with a directive (`try to avoid P') or apprehensional (`beware of P happening') meaning which expands into dependent uses such as negative purpose clauses via extension of the speech reporting construction.
 \newpage
\ea \label{ex:DD80} \textnormal{Evolution of the apprehensive in Mehweb (repeated from \REF{ex:DD59} above)} \\
\ea \label{ex:DD80a} \textnormal{source construction} \\
\textnormal{{\ob}apprehensive{\cb} = warning/directive}
\ex \textnormal{evolution of subordinate uses via extended speech reporting:} \\
\textnormal{{\ob}finite verb {\ob}apprehensive{\cb} \textit{ile}{\cb}  = `fear', negative purpose and `in-case' clauses}
\ex \textnormal{specification of some of the subordinate uses:} \\
\textnormal{{\ob}finite verb {\ob}apprehensive{\cb} \sout{\textit{ile}}{\cb}  = negative purpose and `in-case' clauses}
\z
\z
 
The Akhvakh apprehensive does not have subordinate uses, either because this is a recent development, as suggested by the relative transparency of its connection to the apprehensive interjection, or because it follows a different grammaticalization path. 



\ea \label{ex:DD81} \textnormal{Evolution of the apprehensive in Akhvakh (repeated from \REF{ex:DD78} above)}
\ea \label{ex:DD82a} \textnormal{source construction:} \\
\textnormal{{\ob}verb apprehensive particle{\cb} = warning/directive}
\ex \label{ex:DD83b} \textnormal{univerbation:} \\
\textnormal{{\ob}verb-\textsc{appr}{\cb} = warning/directive}
\z
\z


 A lot remains unclear. We do not know the exact path of grammaticalization of the apprehensive in Akhvakh, and the hypothesis suggested by Denis \citeauthor{Creisselsunpublished2} should be considered with a pinch of salt, probably rather a point of departure for future analyses. In Archi, on syntactic grounds we posit, as a source of the apprehensive, a construction dedicated to `fear' complement clauses, but its morphology remains obscure. Similarly, we know nothing about the morphological origins of the apprehensive in Mehweb, except a hypothetical link to the action nominal in other northern Dargwa varieties. But, again based on syntactic evidence, it seems probable that in Mehweb the form was expanding from main to subordinate clauses, while in Archi, the development was from subordinate to main clauses. And yet, in Archi and Mehweb, the two apprehensives are equally similar in their functional scope to the typological profile of apprehensives cross-linguistically. We conclude this is a case of secondary convergence by filling the same typological niche.  

As we noted above in the discussion of the available data, we have quite a few naturalistic examples from the corpus of Archi, and no textual examples for Mehweb and Akhvakh. All three corpora are narrative, and the apprehensive is not expected to be frequent: as one of the authors of the Akhvakh dictionary and the author of the Akhvakh corpus, Indira Abdulaeva, notices, the Akhvakh apprehensive is after all an \textit{interactional} form. Nevertheless, the absence of the apprehensive in the Mehweb corpus is conspicuous, even if one adjusts for the relative sizes of Archi (over 40,000 tokens) and Mehweb (10,000 tokens) corpora. Indeed, both Archi and Mehweb allow for subordinate uses of the apprehensive such as negative purpose or `fear' complements in elicitation and could be expected to be more proportional in terms of the use of the apprehensive in the narratives. Yet, there are no textual occurrences of the apprehensive in Mehweb, and over 20 naturalistic examples from Archi. This is in line with our hypothesis: the Archi apprehensive has originated from subordinate clause uses in the first place, and is still frequent in subordination, while in Mehweb it has evolved from main clause uses and, in some of its subordinate uses, competes with older alternatives such as purpose converb with a negative suffix.

The three languages of the family where the dedicated apprehensive has so far\footnote{Very late in the process of revising this article, the authors became aware of the presence of forms functionally close to the apprehensive in Tsezic languages. Zaira Khalilova (\citealt{khalilova}) is working on a verbal form \nobreakdash-\textit{dal} in Bezhta, that shows a range of uses similar to those of subordinate uses of the apprehensive in Archi, including `fear' complement and negative purpose. Subordinate uses of the Bezhta form, described as preventative gerund and the negative counterpart of the gerund of aim, are also discussed by \citet[267]{KibrikTestelec2004}. Khalilova does not provide examples of clearly independent, directive uses, definitional for the apprehensive category. The Bezhta form  shares dependent uses with the Archi apprehensive but lacks its independent use, and may correspond to an early stage of the form observed in Archi. More evidence on the origins of this form could thus close the important gap in our argumentation that has not been addressed in the study: what could be the source of a dedicated subordinate `fear' complement form in Archi, which we suggest as the origin of its apprehensive?} been reported belong to three different branches of the family, Lezgic (Archi), Dargwa (Mehweb) and Andic (Akhvakh). Both Archi and Mehweb are outliers of their respective branches in central Daghestan, in strong contact with Avar and Lak but not with each other. The apprehensive does not seem to appear in other Lezgic or Dargwa languages (we have no data for other Andic languages). Akhvakh, an Andic language, also belongs to the area where Avar serves as the main language of interethnic communication (\citealt{DobrushinaZakirova2019}). Areal convergence remains a possibility, but a thorough study of apprehensive constructions in Avar, a major language in contact with all three minority languages discussed above, as well as in Lak, another possible source for two of the three languages, would be required. This could provide a possible missing link crucial for deciding whether areal interaction might have been a (co-)factor in the evolution of the apprehensive forms in the three languages in the study, but is not an easy task because of the dialectal variation within Avar itself, and the fact that the three languages were in contact with different Avar varieties. According to the descriptions of standard varieties, and to our own field data obtained by elicitation, Avar and Lak do not have any dedicated constructions. 


\section{Conclusion}
\label{daniel:conclusion}
 
In this study, we have considered dedicated apprehensives in three East Caucasian languages, all  belonging to different branches of the family and thus only distantly related, Archi (Lezgic), Mehweb (Dargwa) and Northern Akhvakh (Andic). The level of detail of discussion varied according to the amount of available data, and especially to the size of the corpora from which naturalistic examples could be obtained. In all three languages, the source of grammaticalization is obscure, and the path of grammaticalization is probably different. In Akhvakh, the use of the apprehensive is limited to the main clause. For Archi, it is feasible to suggest that the independent apprehensive has developed from `fear' complement clauses via insubordination. For Mehweb, the apprehensive seems to have developed from some kind of a directive form, i.e.\   with main clause uses as a starting point. Yet, following different paths, Archi and Mehweb apprehensives have developed a comparable range of uses, such that they fit the typological range of the apprehensive. All this naturally leads to suggesting an areal convergence scenario. However, none of these languages is in direct contact, and we are unaware of any other East Caucasian language that would have a dedicated morphological apprehensive, including the two potential donors for Archi and Mehweb, namely, Avar and Lak. Apprehensive being a discourse interactional category, one could be inclined to think that areal convergence happens not, or not necessarily, on the morphological level, by developing a dedicated apprehensive form, but on the discourse level, by developing a dedicated construction to express apprehension. Here, Avar, with its dedicated particle \textit{vore(gi)}, could indeed be the attractor for convergence. However, we are unaware of subordinate uses of the apprehensive in Avar that would explain the expansion of the apprehensive to subordinate uses in Mehweb, which led to the remarkable similarity between the two languages. We can only conclude that a more areally informed study of the apprehensives in Daghestanian languages is necessary, and remains a task for further studies. 



\section*{Abbreviations}
\begin{tabularx}{.5\textwidth}{@{}lQ@{}}
\textsc{cl} & agreement slot \\
\textsc{1, 2, 3, 4} & genders in Archi \\
\textsc{m, f, n} & male, female and neuter genders in Mehweb and Akhvakh \\
\textsc{add} & additive particle \\
\textsc{advz} & adverbializer \\
\textsc{all} & allative (goal) \\
\textsc{aor} & aorist \\
\textsc{appr} & apprehensive \\
\textsc{attr} & attributivizer \\
\textsc{aux} & auxiliary \\
\textsc{caus} & causative \\
\textsc{comit} & comitative \\
\textsc{cond} & conditional \\
\textsc{cont} & localization `in contact' \\
\textsc{cop} & copula \\
\textsc{cvb} & converb \\
\textsc{dat} & dative \\
\textsc{dem} & demonstrative \\
\textsc{ego} & egophoric \\
\textsc{erg} & ergative \\
\textsc{ess} & essive (absence of motion as opposed to lative and elative)\\
\textsc{evid} & indirect evidential \\
\textsc{f} & feminine \\
\textsc{f1} & second feminine gender \\
\textsc{fut} & future \\
\textsc{gen} & genitive \\
\textsc{hort} & hortative (first plural inclusive imperative) \\
\textsc{hpl} & human plural (agreement pattern) \\
\textsc{imp} & imperative\end{tabularx}
\begin{tabularx}{.5\textwidth}{@{}lQ@{}}
\textsc{in} & localization `inside'\\
\textsc{incl} & inclusive \\
\textsc{inf} & infinitive \\
\textsc{inter} & localization `between' \\
\textsc{ipfv} & imperfective stem \\
\textsc{irr} & irrealis\\
\textsc{lat} & lative (goal) \\
\textsc{loc} & locative \\
\textsc{locpl} & locutive plural agreement \\
\textsc{lv} & light verb \\
\textsc{n} & neuter gender \\
\textsc{neg} & negative \\
\textsc{nom} & nominative \\
\textsc{npl} & non-human plural (agreement pattern) \\
\textsc{obl} & oblique \\
\textsc{pfv} & perfective stem \\
\textsc{pl} & plural \\
\textsc{proh} & prohibitive (negative imperative) \\
\textsc{prs} & present \\
\textsc{pst} & past \\
\textsc{ptcl} & particle \\
\textsc{pv} & preverb \\
\textsc{purp} & purpose converb\\
\textsc{q} & interrogative particle \\
\textsc{quot} & quotative particle \\
\textsc{self} & pronominal stem used in reflexive and logophoric contexts \\
\textsc{sup} & localization `on' \\
\textsc{temp} & temporal converb \\
\textsc{tr} & transitive \\
V & verb
\end{tabularx}

\section*{Acknowledgements}
The research was partly supported by the joint research funding of the German Federal Government and Federal States in the Academies' Programme, with funding from the Federal Ministry of Education and Research and the Free and Hanseatic City of Hamburg. The Academies’ Programme is coordinated by the Union of the German Academies of Sciences and Humanities. 

We are very grateful to the anonymous reviewers, to Timur Maisak and to the editors of the volume for their important comments on the early version of this paper.


\printbibliography[heading=subbibliography,notkeyword=this]

\end{document}
